%File: anonymous-submission-latex-2026.tex
\documentclass[letterpaper]{article} % DO NOT CHANGE THIS
\usepackage[submission]{aaai2026}  % DO NOT CHANGE THIS
\usepackage{times}  % DO NOT CHANGE THIS
\usepackage{helvet}  % DO NOT CHANGE THIS
\usepackage{courier}  % DO NOT CHANGE THIS
\usepackage[hyphens]{url}  % DO NOT CHANGE THIS
\usepackage{graphicx} % DO NOT CHANGE THIS
\urlstyle{rm} % DO NOT CHANGE THIS
\def\UrlFont{\rm}  % DO NOT CHANGE THIS
\usepackage{natbib}  % DO NOT CHANGE THIS AND DO NOT ADD ANY OPTIONS TO IT
\usepackage{caption} % DO NOT CHANGE THIS AND DO NOT ADD ANY OPTIONS TO IT
\frenchspacing  % DO NOT CHANGE THIS
\setlength{\pdfpagewidth}{8.5in} % DO NOT CHANGE THIS
\setlength{\pdfpageheight}{11in} % DO NOT CHANGE THIS
%
% These are recommended to typeset algorithms but not required. See the subsubsection on algorithms. Remove them if you don't have algorithms in your paper.
\usepackage{algorithm}
\usepackage{algorithmic}
\usepackage{amssymb,amsmath}

% \usepackage{subfigure}
\usepackage{subcaption} % For subfigures


\usepackage{xspace}
% \newcommand{\method}{MetaPrompter\xspace}
\newcommand{\method}{OPTiMACS\xspace}
% \usepackage{verbatimbox}

\usepackage{wrapfig}
\usepackage{hyperref}
\usepackage{booktabs}

%
% These are are recommended to typeset listings but not required. See the subsubsection on listing. Remove this block if you don't have listings in your paper.
\usepackage{newfloat}
\usepackage{listings}
\DeclareCaptionStyle{ruled}{labelfont=normalfont,labelsep=colon,strut=off} % DO NOT CHANGE THIS
\lstset{%
	basicstyle={\footnotesize\ttfamily},% footnotesize acceptable for monospace
	numbers=left,numberstyle=\footnotesize,xleftmargin=2em,% show line numbers, remove this entire line if you don't want the numbers.
	aboveskip=0pt,belowskip=0pt,%
	showstringspaces=false,tabsize=2,breaklines=true}
\floatstyle{ruled}
\newfloat{listing}{tb}{lst}{}
\floatname{listing}{Listing}
%
% Keep the \pdfinfo as shown here. There's no need
% for you to add the /Title and /Author tags.
\pdfinfo{
/TemplateVersion (2026.1)
}

% DISALLOWED PACKAGES
% \usepackage{authblk} -- This package is specifically forbidden
% \usepackage{balance} -- This package is specifically forbidden
% \usepackage{color (if used in text)
% \usepackage{CJK} -- This package is specifically forbidden
% \usepackage{float} -- This package is specifically forbidden
% \usepackage{flushend} -- This package is specifically forbidden
% \usepackage{fontenc} -- This package is specifically forbidden
% \usepackage{fullpage} -- This package is specifically forbidden
% \usepackage{geometry} -- This package is specifically forbidden
% \usepackage{grffile} -- This package is specifically forbidden
% \usepackage{hyperref} -- This package is specifically forbidden
% \usepackage{navigator} -- This package is specifically forbidden
% (or any other package that embeds links such as navigator or hyperref)
% \indentfirst} -- This package is specifically forbidden
% \layout} -- This package is specifically forbidden
% \multicol} -- This package is specifically forbidden
% \nameref} -- This package is specifically forbidden
% \usepackage{savetrees} -- This package is specifically forbidden
% \usepackage{setspace} -- This package is specifically forbidden
% \usepackage{stfloats} -- This package is specifically forbidden
% \usepackage{tabu} -- This package is specifically forbidden
% \usepackage{titlesec} -- This package is specifically forbidden
% \usepackage{tocbibind} -- This package is specifically forbidden
% \usepackage{ulem} -- This package is specifically forbidden
% \usepackage{wrapfig} -- This package is specifically forbidden
% DISALLOWED COMMANDS
% \nocopyright -- Your paper will not be published if you use this command
% \addtolength -- This command may not be used
% \balance -- This command may not be used
% \baselinestretch -- Your paper will not be published if you use this command
% \clearpage -- No page breaks of any kind may be used for the final version of your paper
% \columnsep -- This command may not be used
% \newpage -- No page breaks of any kind may be used for the final version of your paper
% \pagebreak -- No page breaks of any kind may be used for the final version of your paperr
% \pagestyle -- This command may not be used
% \tiny -- This is not an acceptable font size.
% \vspace{- -- No negative value may be used in proximity of a caption, figure, table, section, subsection, subsubsection, or reference
% \vskip{- -- No negative value may be used to alter spacing above or below a caption, figure, table, section, subsection, subsubsection, or reference

%--------------- MACROS ---------
\usepackage{xcolor}
\newcommand{\review}[2]{{\textcolor{#1}{(#2)}}}
\newcommand{\md}[1]{\review{red}{md: #1}}



\setcounter{secnumdepth}{2} %May be changed to 1 or 2 if section numbers are desired.

% The file aaai2026.sty is the style file for AAAI Press
% proceedings, working notes, and technical reports.
%

% Title

% Your title must be in mixed case, not sentence case.
% That means all verbs (including short verbs like be, is, using,and go),
% nouns, adverbs, adjectives should be capitalized, including both words in hyphenated terms, while
% articles, conjunctions, and prepositions are lower case unless they
% directly follow a colon or long dash
% \title{MetaPromptor: RL-based Adaptive Message Structurisation for Multi-Agent Communication}
% \title{Learning Optimal Message Structures for Multi-Agent Communication}
\title{Learning Optimal Message Representations for Multi-Agent Communication}
\author{
    %Authors
    % All authors must be in the same font size and format.
    Written by AAAI Press Staff\textsuperscript{\rm 1}\thanks{With help from the AAAI Publications Committee.}\\
    AAAI Style Contributions by Pater Patel Schneider,
    Sunil Issar,\\
    J. Scott Penberthy,
    George Ferguson,
    Hans Guesgen,
    Francisco Cruz\equalcontrib,
    Marc Pujol-Gonzalez\equalcontrib
}
\affiliations{
    %Afiliations
    \textsuperscript{\rm 1}Association for the Advancement of Artificial Intelligence\\
    % If you have multiple authors and multiple affiliations
    % use superscripts in text and roman font to identify them.
    % For example,

    % Sunil Issar\textsuperscript{\rm 2},
    % J. Scott Penberthy\textsuperscript{\rm 3},
    % George Ferguson\textsuperscript{\rm 4},
    % Hans Guesgen\textsuperscript{\rm 5}
    % Note that the comma should be placed after the superscript

    1101 Pennsylvania Ave, NW Suite 300\\
    Washington, DC 20004 USA\\
    % email address must be in roman text type, not monospace or sans serif
    proceedings-questions@aaai.org
%
% See more examples next
}

%Example, Single Author, ->> remove \iffalse,\fi and place them surrounding AAAI title to use it
\iffalse
\title{My Publication Title --- Single Author}
\author {
    Author Name
}
\affiliations{
    Affiliation\\
    Affiliation Line 2\\
    name@example.com
}
\fi

\iffalse
%Example, Multiple Authors, ->> remove \iffalse,\fi and place them surrounding AAAI title to use it
\title{My Publication Title --- Multiple Authors}
\author {
    % Authors
    First Author Name\textsuperscript{\rm 1},
    Second Author Name\textsuperscript{\rm 2},
    Third Author Name\textsuperscript{\rm 1}
}
\affiliations {
    % Affiliations
    \textsuperscript{\rm 1}Affiliation 1\\
    \textsuperscript{\rm 2}Affiliation 2\\
    firstAuthor@affiliation1.com, secondAuthor@affilation2.com, thirdAuthor@affiliation1.com
}
\fi


% REMOVE THIS: bibentry
% This is only needed to show inline citations in the guidelines document. You should not need it and can safely delete it.
\usepackage{bibentry}
% END REMOVE bibentry

\begin{document}

\maketitle

\begin{abstract}
Large Language Models (LLMs) have demonstrated remarkable capabilities in multi-agent collaborative problem-solving, albeit a gap exists. Existing frameworks predominantly rely on natural language as a primary representation for inter/intra-agentic communication and do not optimize the representation of messages or context that agents communicate/share with one another, which can be indispensable for task-specific performance. 
While recent work has shown that alternative structured representations can enhance reasoning in LLMs, current approaches lack the intelligence necessary to understand, learn or apply optimal communication structures adaptively. 
% We model the problem of dynamically learning the optimal message structures as a Partially Observable Markov Decision Process (POMDP) and propose our method named \method.
We model the problem of dynamically learning the optimal (formal/semi-formal) structures for agentic communication as an Expanding Markov Decision Process (EMDP) and propose our method named \method.
% We introduce \method, a reinforcement learning-based system that adaptively selects and optimizes message formats for multi-agent communication. Our approach models agentic communication as a Partially Observable Markov Decision Process (POMDP) and employs modified off-policy Monte Carlo learning to discover optimal formatting strategies. 
% MetaPromptor seamlessly integrates with existing multi-agent frameworks through a \texttt{modify\_message(source, target, message)} interface, enabling transparent format optimization during agent interactions. 
We evaluate our system across benchmark datasets of collaborative problem-solving with GSMPlus and information exchange via MultiHopQA. 
The results show significant performance improvements while maintaining efficiency.
Our work bridges the gap between rigid communication protocols and open-ended natural language by providing an adaptive framework that learns task-aware structural representations.
% Results demonstrate significant improvements in task completion accuracy while maintaining efficient communication, with our learned formats exhibiting parallels to established agent communication languages. Our work bridges the gap between rigid communication protocols and open-ended natural language by providing an adaptive framework that learns task-specific structural representations through reinforcement learning.
% Large Language Models (LLMs) have demonstrated remarkable capabilities in multi-agent collaborative problem-solving, yet existing frameworks predominantly rely on natural language communication without optimizing message formats for task-specific performance. While recent work has shown that alternative structured formats can enhance single-LLM reasoning and reduce token usage in multi-agent settings, current approaches lack adaptive mechanisms to learn optimal communication structures dynamically. We introduce MetaPromptor, a reinforcement learning-based system that adaptively selects and optimizes message formats for multi-agent communication. Our approach models agentic communication as a Partially Observable Markov Decision Process (POMDP) and employs modified off-policy Monte Carlo learning to discover optimal formatting strategies. MetaPromptor seamlessly integrates with existing multi-agent frameworks through a \texttt{modify\_message(source, target, message)} interface, enabling transparent format optimization during agent interactions. We evaluate our system across three domains: collaborative problem-solving with GSMPlus mathematical reasoning, information exchange via MultiHopQA datasets (NarrativeQA, WikiHopQA, HotpotQA), and the FLASH multi-agent triage framework. Results demonstrate significant improvements in task completion accuracy while maintaining efficient communication, with our learned formats exhibiting parallels to established agent communication languages. Our work bridges the gap between rigid communication protocols and open-ended natural language by providing an adaptive framework that learns task-specific structural representations through reinforcement learning.


\end{abstract}

% Uncomment the following to link to your code, datasets, an extended version or similar.
% You must keep this block between (not within) the abstract and the main body of the paper.
% \begin{links}
%     \link{Code}{https://aaai.org/example/code}
%     \link{Datasets}{https://aaai.org/example/datasets}
%     \link{Extended version}{https://aaai.org/example/extended-version}
% \end{links}

\section{Introduction}\label{intro}

% Give context of communication and different forms and motivate the question that is NL the best form
Communication is essential among humans and forms the very fabric of society. Along with natural language, we convey information using gestures, mathematical expressions, multimedia etc. 
% \md{Is the above sentence appropriate - with this argument reviewers would think we are approaching the multi-modal space}. 
Depending on the task, different forms of communication are important to accomplish the end goal in an effective and efficient manner. For example, in a driving scenario, the vehicles and the traffic system efficiently coordinate primarily by visual means using colored lights, signs, markers etc. rather than relying solely on natural language (NL).  In the realm of human-computer interaction, traditional forms of communication primarily consisted of user commands (using text, touch, gestures) that perform desired action through APIs encoded in high or low level machine instructions. Lately, with the advent of Large Language Models (LLMs) natural language has become the dominant mode of communication, between LLM and human agents and between LLM agents. This prevalence extends to multi-agent systems, where specialized agents collaborate through natural language exchanges to solve complex tasks \cite{li2023camel,wu2023autogen,park2023generative}. But one question naturally arises, even among text based agents, whether natural language in general is the best representation of communication to solve a given task of varying complexity and nature (cf. fig. \ref{fig:motivation})? 
% \md{Modality is a whole new argument - to highlight the pitfalls of NL we do not essentially need to spend a paragraph on other modalities - instead we need to highlight the fuzzyness and ambiguities of NL - the lack formalisms, lack of consistency and the fact that language used by humans is subject to a lot unsaid, implicitly understood subtext and context}. 
% Natural language has long served as the primary medium for human-computer interaction and remains the dominant communication format in contemporary Large Language Model (LLM) applications. This prevalence extends to multi-agent systems, where specialized agents collaborate through natural language exchanges to solve complex tasks \cite{li2023camel,wu2023autogen,park2023generative}. However, the inherent ambiguities and inefficiencies of natural language may not align optimally with the precision requirements and computational constraints of multi-agent coordination.

\begin{figure}[t!]
    \centering
    \includegraphics[width=0.99\linewidth]{AAAI/figures/sac_motivation.png}
    \caption{In agentic systems, optimal message representations (JSON, proposition, code etc.) enable agents to collaborate more effectively, reduce \emph{ambiguity, bias}, prevent error propagation and improve task performance.}
    \label{fig:motivation}
\end{figure}

% Talk about recent works that show there are other forms of communication than NL
Recent advances in LLM reasoning have challenged the assumption that natural language represents the optimal intermediate representation for cognitive processes. 
Many works such as Program-of-Thought \cite{chen2022program}, X-of-Thought \cite{ding2024thoughtsdefyinglawpenrose}, PLaG \cite{plag_icml} have demonstrated that alternative structured representations can enhance reasoning and decision making capabilities. This ineffectiveness of using NL is compounded in multi-agent settings where multiple LLM agents are involved and pass messages using natural language as the de facto mode. In this vein, recent research \cite{chen-etal-2024-beyond-natural} suggests that employing non-natural language as intermediate formats can help achieve improvements in reasoning performance during agentic communication.
% efficiency and up to 72.7\% reductions in token usage during multi-agent communication.

% Talk about limitations in current works
Despite the promising developments in literature, current single and multi-agent frameworks lack a principled and data driven approach to adaptively select the optimal format of the intermediate representation based on task requirements. Existing approaches rely on static format selection or on hand crafted prompts that let the LLM choose the appropriate format for the task based on its inductive biases. 
% The former is based on the premise that the specified communication representation is the best for the given task whereas the latter suffers from the inductive biases of the LLM. 
Both the static formats and the ones suggested by the LLM may not be optimal. 
Moreover, the prior works rely on a fixed set of structures, whereas new optimal representations could dynamically emerge  as the agent interacts with the environment. The fundamental challenge lies in developing systems that can \emph{learn} optimal message formatting strategies \emph{dynamically} from historical \emph{data} or from the \emph{environment} while maintaining \emph{seamless integration} with existing systems.
% multi-agent architectures. 
% Moreover in multi-agent settings, natural language is the de facto choice of communication between the LLM agents. 

% Despite these promising developments, current multi-agent frameworks lack principled mechanisms for adaptively optimizing communication formats based on task requirements and agent capabilities. Existing approaches either rely on static format selection or require manual specification of communication structures, limiting their applicability across diverse problem domains. The fundamental challenge lies in developing systems that can learn optimal message formatting strategies dynamically while maintaining seamless integration with existing multi-agent architectures.

% Our paper and how it differs
In this paper, we address the challenge of learning the optimal representation to solve a given sub task of a problem. We propose to learn the optimal representations in a \emph{data-driven} manner and from the observations/rewards as the agents interact with the environment. Since novel optimal representations could emerge as the agent learns, we further propose to dynamically adapt to the set of tasks and communication representations rather than optimizing for a fixed set of representations. 

% Motivation and Some details
% We propose to \emph{learn} optimal structures (action) \emph{dynamically} for a given task (state), as the agents interact with the environment. So we naturally model the problem as a Markov Decision Process (MDP).
% In order to \emph{learn} optimal representations (action) \emph{dynamically} for a given task (state) as the agents interact with the environment, we naturally model the problem as a Markov Decision Process (MDP).
We aim to \emph{learn} optimal representations (action) \emph{dynamically} for a given task (state) as the agents interact with the environment, \textcolor{blue}{e.g. what is the best format for an equation solving task. This naturally reduces to a sequential decision-making problem which can be modeled as a Markov Decision Process (MDP).
As the representations could evolve dynamically, we work in a setting where the states and actions are evolving/expanding. Thus, we pose the optimal representation selection problem as an Expanding
% /Expanding 
Markov Decision Process (EMDP) and introduce our method called \method - \emph{\underline{O}ptimal \underline{P}rompt \underline{T}ransformation \underline{i}n \underline{M}ulti-\underline{A}gent \underline{C}ommunication \underline{S}ystem}.} Considering the agents' communication trajectory of length $L_T$ and our desire to learn the optimal message representation at each step/turn, a naive approach to learn the optimal representation for a task from the historical data would incur a complexity of $\mathcal{O}((\lvert\mathcal{T}\lvert\lvert\mathcal{F}\lvert)^{L_T})$, where $\mathcal{T}, \mathcal{F}$ (cf. $\S$ \ref{computational_complexity}) are the set of tasks and representations/formats respectively. We thus employ methods from RL to efficiently learn the state action space in linear complexity. 

\begin{figure}[h!]
    \centering
    \includegraphics[width=0.70\linewidth]{AAAI/figures/sac_single_agentV2.png}
    \caption{Task performance improves when prompt is formatted as per discovered optimal representation.}
    \label{fig:single_agent}
\end{figure}

Motivating the need to learn the optimal representations in a data driven manner, consider the following problem assigned to the system (cf. fig. \ref{fig:system_overview}): ``Gerald works at a daycare that pays him \$30 every day. He worked for an \emph{entire week},...he spent \$100 on... How much money does he have left? ''. The problem states the person works for 7 days and saves \$110. Plain text (NL) based solution arrives at the answer of \$50, which stems from the NL bias that a working week consists of 5 days. 
% Our method selects code as the message representation and correctly sets the variable of number of working days to be a week (=7 days) and arrives at the expected answer. 
Prompting the LLM to think in code (suggested by our method) enables the model to explicitly define the number of days as 7, thereby eliminating \emph{common biases associated with NL reasoning}, arriving at the expected answer. 
Further examples are in the \textcolor{blue}{Appendix $\S$ \ref{case_study}}. To make the motivation concrete, we perform a study on single agent (LLM) tasks where we compare results of plain text, optimal representation decided by LLM and optimal representation derived by exploring the data (our method). We present the findings in fig. \ref{fig:single_agent} where we see that our method consistently outperforms the baselines (p-value<0.05 using Student's t-test, std. dev. < 2\%, LLM: GPT-4). The errors induced by natural language ambiguity get amplified in a multi-agent scenario with many turns, necessitating the discovery of task specific optimal message representations to input to the LLM agent.



% Contributions
\noindent \textbf{Contributions: } To summarize, we make the following contributions: (1) We formalize the problem of selecting the optimal communication representation in a multi-agent system as an evolving MDP (2) We propose our method \method, to learn the optimal policy for task and message representation in a \emph{data driven} and \emph{dynamic manner}. (3) Extensive experiments are conducted on 7 benchmark datasets to show the efficacy of our method. 
For reproducibility we provide pseudo-code and prompts in Appendix (and code post acceptance).
% The code will be made public post acceptance (cf. blind review).

% We address this challenge by introducing MetaPromptor, a reinforcement learning-based framework that treats message format selection as a sequential decision-making problem. Our key insight is that different LLMs exhibit varying sensitivities to structural representations, making adaptive format selection crucial for optimal performance, particularly in resource-constrained environments with smaller language models. Furthermore, structured formats enable formal reasoning about artifacts and knowledge while reducing representational ambiguity \cite{li2023structured}.

% MetaPromptor models multi-agent communication as a Partially Observable Markov Decision Process (POMDP), where agents must select optimal message formats based on partially observable conversation states and task contexts. We approximate this POMDP as a standard MDP when the state-action space evolves slowly relative to the learning timescale, enabling tractable Q-learning while handling residual partial observability through cyclic expansion-convergence strategies.

% Our system operates through three core components: (1) task classification for state identification, (2) hybrid format selection combining learned Q-values with LLM-generated formats, and (3) off-policy Monte Carlo learning with dynamic action space expansion. The format selection mechanism integrates multiple strategies: exploitation of learned Q-values through softmax sampling, LLM-guided generation considering historical performance, and exploration through inverse frequency sampling. This multi-faceted approach ensures both exploitation of proven formats and discovery of novel structures.

% We evaluate MetaPromptor across three distinct domains: collaborative problem-solving using GSMPlus mathematical reasoning benchmarks, information exchange through MultiHopQA datasets (NarrativeQA, WikiHopQA, HotpotQA), and the FLASH multi-agent triage framework for incident prioritization. Our results demonstrate consistent improvements in task completion accuracy while maintaining communication efficiency.

% The contributions of this work are threefold: (1) We formalize multi-agent communication format optimization as a reinforcement learning problem, providing a principled framework for adaptive message structuring; (2) We introduce MetaPromptor, a practical system that seamlessly integrates with existing multi-agent frameworks while learning optimal communication formats; (3) We demonstrate empirical improvements across diverse domains, revealing that learned formats naturally converge toward established agent communication paradigms.

\section{Related Works}

Here we briefly compare our work with related works for LLM agents and defer detailed discussion to the appendix.

% \subsection{Multi-Agent Systems (MAS)}
% Recent advances in MAS have evolved from single-shot LLM interactions to structured communication \cite{zhao2023survey}. The framework proposed in \cite{googadk2025} formalizes the model-orchestration-tools stack with persistent memory systems.
% % while their Agent Development Kit (ADK) supports complex multi-agent orchestration with persistent memory systems \cite{google2023adk}. 
% \cite{magenticone2024} demonstrates practical coordination through specialized agents 
% % (WebSurfer, FileSurfer, Coder, ComputerTerminal) 
% managed by a central Orchestrator. 
% % These systems exemplify four core agentic reasoning patterns: reflection \cite{madaan2023self,shinn2023reflexion}, tool-use, planning, and multi-agent collaboration \cite{wang2023survey}. 
% However, existing frameworks focus primarily on agent coordination mechanisms rather than optimizing communication structures between agents.

% \subsection{Multi-Agent Communication Protocols}
\subsection{Agentic Communication Protocols}
Traditional agent communication has relied on standardized protocols like KQML \cite{finin1994kqml}, FIPA-ACL \cite{fipa1997acl}, MASIF \cite{milojicic1999masif}, Web Services \cite{berners2001web}, Event Stream Processing \cite{cugola2012processing} etc.
In the era of LLMs researchers have explored communication using function calling mechanisms \cite{schick2023toolformer} and tool usage \cite{schick2023toolformer}.
% Traditional agent communication has relied on standardized protocols like KQML \cite{finin1994kqml} and FIPA-ACL \cite{fipa1997acl}, with MASIF \cite{milojicic1999masif} enabling cross-platform mobility. Modern approaches integrate Web Services \cite{berners2001web}, Event Stream Processing \cite{cugola2012processing}, and RAG architectures \cite{lewis2020retrieval} for enhanced grounding capabilities. Function calling mechanisms \cite{schick2023toolformer} and Toolformer \cite{schick2023toolformer} demonstrate autonomous tool usage through structured interfaces.
Recent standardization efforts include protocols such as MCP \cite{anthropic2024mcp}, ACP \cite{acp2024protocol}, A2A \cite{a2a2024protocol} and ANP \cite{anp2024protocol}. While these establish agent communication infrastructure, they provide static, rule-based schemas that cannot adapt to varying task requirements or agent capabilities.
% Recent standardization efforts include the Model Context Protocol (MCP) \cite{anthropic2024mcp} for LLM context ingestion, Agent Communication Protocol (ACP) \cite{acp2024protocol} for REST-native messaging, Agent-to-Agent Protocol (A2A) \cite{a2a2024protocol} for capability-based delegation, and Agent Network Protocol (ANP) \cite{anp2024protocol} for decentralized agent discovery. While these protocols establish communication infrastructure, they provide static, rule-based schemas that cannot adapt to varying task requirements or agent capabilities.

% \subsection{Learning in Multi-Agent Systems}
% Multi-Agent Reinforcement Learning (RL) has explored communication optimization through various approaches. 
% Message pruning \cite{singh2019message} and graph based \cite{zhang2019graph} methods focus on reducing communication overhead and optimizing information flow between agents. 
% % Message pruning techniques \cite{singh2019message} focus on reducing communication overhead while maintaining coordination effectiveness. Graph-based coordination methods \cite{zhang2019graph} leverage network structures to optimize information flow between agents. 
% % Emergent symbolic languages \cite{havrylov2017emergence} demonstrate how agents can develop communication protocols through interaction.
% Task-oriented communication protocols have been developed using RL to enable agents to discover efficient communication strategies between classical RL agents \cite{karten2023role,foerster2016learning}. However LLMs differ in linguistically mediated communication \cite{sumers2023cognitive}.
% Existing multi-agent LLM frameworks leverage world knowledge for coordination through open-loop and closed-loop methods \cite{shinn2023reflexion}, but do not adapt message structures to optimize task performance.
% % However, these works primarily focus on classical RL agets and Language Model agents exhibit fundamental differences, including statelessness, stochastic behavior, semantic sensitivity, and linguistically mediated communication \cite{sumers2023cognitive}.
% % While LM agent scaffolding approaches employ extended memory systems, tool integration, and cognitive architectures \cite{sumers2023cognitive}, 
% % the communication optimization problem remains largely unaddressed. 

\subsection{Structured Communication among LLM Agents}

\cite{yuan2023structured} prompt LLMs to generate optimal formats for fine-tuning tasks, but do not explore inter-agent communication contexts.
Recent work on structured communication focused on leveraging the inductive biases in LLM for format generation \cite{chen2022program}, though without mechanisms for dynamically learning or domain adaptation. These methods suffer from over-reliance on static LLM biases without learning from environment feedback.
% adaptive learning mechanisms.
% The AutoForm framework \cite{chen-etal-2024-beyond-natural} demonstrates that LLMs can autonomously select non-natural language formats for single-LLM reasoning and multi-agent communication, achieving 3.3-5.7\% improvements in reasoning efficiency and up to 72.7\% reductions in token usage. This work reveals that LLMs can devise suitable formats from task-specific examples and that these formats are transferable across different LLMs. Importantly, the communication formats adopted by LLMs mirror traditional Agent Communication Languages like KQML, highlighting their clarity, brevity, and structured nature.
% Parallel work in chain-of-thought reasoning demonstrates that reasoning quality depends more on structural patterns than factual content \cite{wang2023plan}. We extend this insight to multi-agent communication: message effectiveness depends crucially on structural encoding rather than content alone, suggesting that optimal communication strategies emerge from adaptive format selection rather than fixed schemas.

Our work differs from existing research in that, unlike existing approaches that focus on either agent coordination infrastructure or static / LLM prompted format generation, we focus on the optimization of message representations/structures/formats in multi-agent communication by \emph{dynamically learning} from the environment.

% % \section{Changes}
% % \subsection{Related Work}
% [1] Add intro from Optima + Justification of the 2 types of setting : https://arxiv.org/pdf/2410.08115
% [2] The Prompts are inspired from FIPA-KQQML ACLs. Though we could select any states, we chose the task as the states based on Autoform (2-level)
% [3] MultiHopQA DataSet Preprocessing inspired by A2 Appendix in Autoform
% [4] In Related Work, add AgentPrune and Optima works.
% [5] For Debate, cite Due et al, ChatEval (again from Optima Paper)
% [6] For 0.7 as threshold, we picked this from works: \cite{gokhan2024regnlp, louis2023opinesum, shen2023quantifying} 
% [7] Justify the DataPoint Numbers from Autoform
% [8] Why need of dynamic structuriser - (1) dynamism helps break inductive biases, where for a given task, we start with some bootstrapped actions and increment the set of actions through LLM's inductive bias; while scoring them (as Q-values). The need for dynamic structurizers stems from their ability to overcome the fundamental limitation that asking LLMs to directly generate the optimal action is less effective than asking them to generate multiple candidate actions and selecting among them. This principle is well-established in the pass@k literature, where generating k candidate solutions and selecting the best one consistently outperforms single-shot generation (pass@1) \cite{chen2021evaluating,wang2022self}. This superiority of multiple candidate generation extends naturally to action selection in reinforcement learning settings. Best-of-N (BoN) sampling strategies demonstrate that generating multiple candidate outputs and selecting the optimal one based on a scoring function significantly outperforms single-output generation \cite{ichihara2025evaluation,jinnai2024regularized}. In dynamic action spaces, this becomes even more critical, as "repeated sampling with multiple LLMs yield better results than with a single LLM" because different models can generate complementary candidate actions that collectively cover a broader solution space \cite{chen2024truly}. Dynamic structurizers leverage this principle by continuously generating diverse action candidates through LLM's inductive bias, then scoring them as Q-values to identify optimal actions. This approach enables the system to break free from static action biases by maintaining a dynamic exploration of the action space, where "a prior optimal policy plays an important role in action pick-up by providing useful knowledge and experience" when new actions emerge \cite{dang2023action}. The theoretical foundation for this improvement lies in the fact that multiple candidate generation provides better coverage of the solution space, allowing the system to discover actions that a single greedy generation would miss.
% (2) dynamism makes the structuriser robust to data-shifts in an online-setting where it could start with prior beliefs about the values, and gradually update them while discovering the new beliefs.
% [9] Analyse structures based on : (1) training and model architecture - Paper : Remnants of Autoregression; Style Transfer Metrics (2) Prompt Optimisation Methods (Literature Review)
% [10] 



% Recent advances in multi-agent systems have evolved from single-shot LLM interactions to sophisticated autonomous architectures capable of reasoning, planning, and coordination through structured communication \cite{zhao2023survey}. Google's three-layer framework formalizes the model-orchestration-tools stack \cite{google2023framework}, while their Agent Development Kit (ADK) supports complex multi-agent orchestration with persistent memory systems \cite{google2023adk}. Microsoft's Magnetic-One demonstrates practical coordination through specialized agents (WebSurfer, FileSurfer, Coder, ComputerTerminal) managed by a central Orchestrator \cite{microsoft2023magnetic}.

% These systems exemplify four core agentic reasoning patterns: reflection \cite{madaan2023self,shinn2023reflexion}, tool-use, planning, and multi-agent collaboration \cite{wang2023survey}. However, existing frameworks focus primarily on agent coordination mechanisms rather than optimizing communication formats between agents.

% Traditional agent communication has relied on standardized protocols like KQML \cite{finin1994kqml} and FIPA-ACL \cite{fipa1997acl} for speech acts, with MASIF \cite{milojicic1999masif} enabling cross-platform mobility. Modern approaches integrate Web Services \cite{berners2001web}, Event Stream Processing \cite{cugola2012processing}, and RAG architectures \cite{lewis2020retrieval} for enhanced grounding capabilities. Function calling mechanisms \cite{schick2023toolformer} and Toolformer \cite{schick2023toolformer} demonstrate autonomous tool usage through structured interfaces.

% Recent standardization efforts include the Model Context Protocol (MCP) \cite{anthropic2024mcp} for LLM context ingestion, Agent Communication Protocol (ACP) \cite{acp2024protocol} for REST-native messaging, Agent-to-Agent Protocol (A2A) \cite{a2a2024protocol} for capability-based delegation, and Agent Network Protocol (ANP) \cite{anp2024protocol} for decentralized agent discovery. While these protocols establish communication infrastructure, they provide static, rule-based schemas that cannot adapt to varying task requirements or agent capabilities.

% \subsection{Multi-Agent Reinforcement Learning for Communication}

% Multi-Agent Reinforcement Learning has explored communication optimization through various approaches. Message pruning techniques \cite{singh2019message} focus on reducing communication overhead while maintaining coordination effectiveness. Graph-based coordination methods \cite{zhang2019graph} leverage network structures to optimize information flow between agents. Emergent symbolic languages \cite{havrylov2017emergence} demonstrate how agents can develop communication protocols through interaction.

% Task-oriented communication protocols have been developed using reinforcement learning to enable agents to discover efficient communication strategies \cite{karten2023role,foerster2016learning}. However, Language Model agents exhibit fundamental differences from classical RL agents, including statelessness, stochastic behavior, semantic sensitivity, and linguistically mediated communication \cite{sumers2023cognitive}.

% While LM agent scaffolding approaches employ extended memory systems, tool integration, and cognitive architectures \cite{sumers2023cognitive}, the communication optimization problem remains largely unaddressed. Existing multi-agent LLM frameworks leverage world knowledge for coordination through open-loop (ReAct, ADaPT) and closed-loop methods (Refiner, Retroformer, REX) \cite{shinn2023reflexion}, but do not adapt message formats to optimize task performance.

% \subsection{Structured Communication and Format Optimization}

% Early work on structured communication focused on leveraging LLM inductive biases for format generation \cite{chen2022program}, though without mechanisms for online learning or domain adaptation. Recent approaches prompt LLMs to generate optimal formats for fine-tuning tasks \cite{yuan2023structured}, but do not explore inter-agent communication contexts. These methods suffer from over-reliance on static LLM biases without adaptive learning mechanisms.

% The AutoForm framework \cite{chen-etal-2024-beyond-natural} demonstrates that LLMs can autonomously select non-natural language formats for single-LLM reasoning and multi-agent communication, achieving 3.3-5.7\% improvements in reasoning efficiency and up to 72.7\% reductions in token usage. This work reveals that LLMs can devise suitable formats from task-specific examples and that these formats are transferable across different LLMs. Importantly, the communication formats adopted by LLMs mirror traditional Agent Communication Languages like KQML, highlighting their clarity, brevity, and structured nature.

% Parallel work in chain-of-thought reasoning demonstrates that reasoning quality depends more on structural patterns than factual content \cite{wang2023plan}. We extend this insight to multi-agent communication: message effectiveness depends crucially on structural encoding rather than content alone, suggesting that optimal communication strategies emerge from adaptive format selection rather than fixed schemas.

% \subsection{Neurosymbolic Systems and Representation Learning}

% Contemporary neurosymbolic systems insert explicit, machine-verifiable representations between language generation and reasoning. SymbCoT translates natural language problems into first-order logic for symbolic solving \cite{zhang2023symbcot}, while VERUS-LM decouples knowledge construction from inference using satisfiable logical theories \cite{bogdan2023verus}. These approaches optimize solver fidelity while maintaining fixed intermediate schemas.

% Program-of-Thought (PoT) \cite{chen2022program,gao2023pal} prompts models to generate code as thought and offloads answer generation to code interpreters. X-of-Thought \cite{liu2023x} integrates CoT, PoT, and Equation-of-Thought, dynamically ensembling these methods for improved reasoning. Tree-of-Thought \cite{yao2023tree} employs depth and breadth-first search techniques to produce high-quality reasoning chains.

% While some CoT variants explore formats beyond natural language for reasoning, the chosen formats' improvements are often obscured by concurrent use of supplementary tools such as code interpreters, blurring the distinction between format efficacy and tool execution.

% \subsection{Our Contribution}

% Unlike existing approaches that focus on either agent coordination infrastructure or static format generation, MetaPromptor addresses the adaptive optimization of message formats in multi-agent communication. Our system bridges the gap between rigid communication protocols and open-ended natural language by learning task-specific structural representations through reinforcement learning.

% MetaPromptor diverges from neurosymbolic approaches by treating the schema itself as a learnable component, using reinforcement learning to select and evolve message formats based on downstream task success. This creates a structure-aware system that remains adaptively plastic—capable of tailoring representations to specific agent pairs, task contexts, and performance constraints that static symbolic formalisms cannot accommodate.

% By treating format selection as a POMDP and employing off-policy Monte Carlo learning with dynamic action space expansion, MetaPromptor can discover and optimize communication formats that enhance multi-agent task performance—a capability not addressed by current agent frameworks or communication protocols.

% \section{Prelimnaries}
% \section{Preliminaries and Problem Formulation}
\section{Problem Formulation}

% In this section, we elucidate certain concepts key to understanding our method and layout the problem.
In this section, we formulate the problem setup. For readers unaware of Markov Decision Process (MDP) and related techniques we kindly refer to the \textcolor{blue}{Appendix (\ref{prelims})} or \cite{sutton2018reinforcement}.

% \subsection{Preliminaries}

% \noindent \textbf{Markov Decision Process (MDP)} provides a mathematical framework for modeling decision-making assuming \textit{Markov property}, meaning the future state depends only on the current state and action. Formally, an MDP is defined as a tuple $(\mathcal{S}, \mathcal{A}, \mathcal{P}, \mathcal{R}, \gamma)$, where $\mathcal{S}, \mathcal{A}$ are the set of states \& actions, $\mathcal{P}(s'|s,a)$ is the transition probability, $\mathcal{R}(s,a)$ is the reward function and $\gamma \in [0,1]$ is the discount factor.
% The objective in an MDP is to find a policy $\pi: \mathcal{S} \rightarrow \mathcal{A}$ that maximizes the expected cumulative discounted reward, defined as: $V^\pi(s) = \mathbb{E}_\pi \left[ \sum_{t=0}^{\infty} \gamma^t R(s_t, a_t) \mid s_0 = s \right]$

% \noindent \textbf{Partially Observable Markov Decision Processes (POMDPs)} provide a principled framework for sequential decision-making under uncertainty, where the true state of the environment is not directly observable. In a POMDP, the agent maintains a belief (a probability distribution) over possible states, updating this belief as it receives observations and takes actions. Formally, a POMDP is defined by the tuple $(\mathcal{S}, \mathcal{A}, \mathcal{O}, \mathcal{P}, \mathcal{R}, \mathcal{\ohm}, \gamma)$, where $\mathcal{O}$ is the set of observations and $\gamma$ is the discount factor. POMDPs are known to be intractable in practice, especially in multi-agent or high-dimensional settings. The core challenge lies in the need to reason over the entire belief space, which grows exponentially with the number of states and agents. Exact solutions require maintaining and updating a belief over all possible states at every timestep, leading to computational complexity that is prohibitive for real-world problems, particularly when the state and action spaces are large or evolving.

% % \noindent \textbf{Q-Learning} is a model-free RL algorithm that aims to learn the optimal action-selection policy by estimating the action-value function, denoted as $Q(s,a)$. The core idea of Q-learning is to iteratively update the Q-values using the Bellman equation: $Q(s_t, a_t) \leftarrow Q(s_t, a_t) + \alpha \left[ r_t + \gamma \max_{a'} Q(s_{t+1}, a') - Q(s_t, a_t) \right]$.

% \noindent \textbf{Monte Carlo Learning} method for learning Q-values, is a model-free approach that estimates the action-value function $Q(s,a)$ using complete episodes of experience. For each state-action pair $(s,a)$ encountered in an episode, the Q-value is updated as: $Q(s,a) \leftarrow Q(s,a) + \alpha \left[ G_t - Q(s,a) \right]$
% where $\alpha$ is the learning rate and $G_t = \sum_{k=t}^{T} \gamma^{k-t} r_k$ is the return from time step $t$ until the end of the episode $T$. 
% For further details on RL we refer the reader to ~\cite{sutton2018reinforcement}.


% % \subsubsection{Markov Decision Process (MDP)}

% % A Markov Decision Process (MDP) provides a mathematical framework for modeling decision-making in environments where outcomes are partly random and partly under the control of an agent. Formally, an MDP is defined as a tuple $(\mathcal{S}, \mathcal{A}, P, R, \gamma)$, where:
% % \begin{itemize}
% %     \item $\mathcal{S}$ is the set of states,
% %     \item $\mathcal{A}$ is the set of actions,
% %     \item $P(s'|s,a)$ is the transition probability function, representing the probability of moving to state $s'$ from state $s$ after taking action $a$,
% %     \item $R(s,a)$ is the reward function, giving the expected immediate reward received after taking action $a$ in state $s$,
% %     \item $\gamma \in [0,1]$ is the discount factor, which balances the importance of immediate and future rewards.
% % \end{itemize}

% % The objective in an MDP is to find a policy $\pi: \mathcal{S} \rightarrow \mathcal{A}$ that maximizes the expected cumulative discounted reward, defined as:
% % \[
% % V^\pi(s) = \mathbb{E}_\pi \left[ \sum_{t=0}^{\infty} \gamma^t R(s_t, a_t) \mid s_0 = s \right],
% % \]
% % where $V^\pi(s)$ is the value function under policy $\pi$. MDPs assume the \textit{Markov property}, meaning the future state depends only on the current state and action, not on the sequence of events that preceded it. This framework underpins many reinforcement learning algorithms and is foundational to understanding agent-environment interactions~\cite{puterman1994markov}.

% % \subsubsection{Policy Learning (Q-learning)}
% % \textbf{Q-learning} is a model-free reinforcement learning algorithm that aims to learn the optimal action-selection policy by estimating the action-value function, denoted as $Q(s,a)$. The action-value function represents the expected cumulative discounted reward of taking action $a$ in state $s$ and following the optimal policy thereafter. The core idea of Q-learning is to iteratively update the Q-values using the Bellman equation:
% % \[
% % Q(s_t, a_t) \leftarrow Q(s_t, a_t) + \alpha \left[ r_t + \gamma \max_{a'} Q(s_{t+1}, a') - Q(s_t, a_t) \right],
% % \]
% % where $\alpha \in (0,1]$ is the learning rate, $\gamma \in [0,1]$ is the discount factor, $r_t$ is the reward received after taking action $a_t$ in state $s_t$, and $s_{t+1}$ is the resulting next state. Over time, and under certain conditions such as sufficient exploration and decaying learning rate, Q-learning converges to the optimal action-value function $Q^*(s,a)$, enabling the derivation of the optimal policy $\pi^*(s) = \arg\max_a Q^*(s,a)$~\cite{watkins1992qlearning}.

% % \subsubsection{Monte Carlo Learning}
% % The Monte Carlo (MC) method for Q-learning, is a model-free approach that estimates the action-value function $Q(s,a)$ using complete episodes of experience. Unlike temporal-difference methods, Monte Carlo methods update Q-values only at the end of an episode, using the actual returns observed. For each state-action pair $(s,a)$ encountered in an episode, the Q-value is updated as:
% % \[
% % Q(s,a) \leftarrow Q(s,a) + \alpha \left[ G_t - Q(s,a) \right],
% % \]
% % where $\alpha$ is the learning rate and $G_t = \sum_{k=t}^{T} \gamma^{k-t} r_k$ is the return from time step $t$ until the end of the episode $T$. Monte Carlo methods require episodes to terminate and assume that the environment is episodic. They are particularly useful when the model dynamics are unknown and when bootstrapping is not desired. While slower to converge than TD methods, MC-based Q-learning can provide unbiased estimates of the expected return under sufficient exploration~\cite{sutton2018reinforcement}.


% \subsection{Problem Formulation}

We operate in a multi-agentic setting where multiple specialized LLM agents collaborate \emph{sequentially} to solve a complex task. 
% As such effective communication among such agents over multiple turns is a key driver of a performant solution. 
Let there be $N$ LLM agents $\{\mathfrak{A}_i\}_{i=1}^N$, passing messages to each other over $T$ turns. In any turn $t, 1 \leq t \leq T$, agent $\mathfrak{A}_i$ sends a message $m_{(i,j)}^t$ to next agent $\mathfrak{A}_j$ in the pipeline to share information and context. We aim to develop a 
% functional transformation 
module $F(m_{(i,j)}^t) \rightarrow m'$ that transforms message from $\mathfrak{A}_i$ to $\mathfrak{A}_j$. We optimize $m'=F(m_{(i,j)}^t)$ such that $m'$ is in a (near) optimal structure/representation (defined by goal success) given the complex task and the nature of $\mathfrak{A}_i$ and $\mathfrak{A}_j$. 
% In particular, we aim to learn a dynamic extensible policy $\pi^*(F(m)|\text{Task}_i,m,t,\mathfrak{A}_i, \mathfrak{A}_j)$
Considering we desire to learn the optimal structures (action) for a given task (state), we naturally pose the problem as a Markov decision process (MDP) \cite{puterman1994markov} and aim to learn the optimal policy $\pi^*(F|\text{Task}_i,m,t,\mathfrak{A}_i, \mathfrak{A}_j)$ which maps a given task to the optimal representation.
\textcolor{blue}{Simpler alternatives like contextual bandits \footnote{\textcolor{blue}{In the single agent (1-step) setup the MDP simplifies to a Contextual Bandit. Refer Appendix \ref{sec:mdp_justification} for details.}} would not solve for the sequential nature of the problem and A/B testing to find the best format is not practical (cf. $\S$ \ref{sec:mdp_justification}).}
Formally, we aim to learn the policy $\pi^*$ such that the end goal of the multi-agent communication is achieved. This can be represented as:
% \begin{equation}
%     \pi^*(F(m)|\text{Task}_i,m,t,\mathfrak{A}_i, \mathfrak{A}_j) = \underset{\pi \in \Pi}{\text{argmax}} \quad f_s(\text{Task}_j)
% \end{equation}
% \small{
% \begin{equation}
%     \pi^* = \underset{\pi \in \Pi}{\text{argmax}} \quad f_s(\text{Task}_j, F(m) | f = \arg\max_{f\in\mathcal{F}}\pi(F|\text{Task}_i,m,t,\mathfrak{A}_i, \mathfrak{A}_j))
% \end{equation}
% }
% % \small{
% \begin{equation}
%     \begin{aligned}
%         \pi^* = \underset{\pi \in \Pi}{\text{argmax}} \sum_{\text{task}} & f_s(\text{task}, F(m) | \\
%         &F = \arg\max_{f\in\mathcal{F}}\pi(f|\text{context})) 
%     \end{aligned}
% \end{equation}
% % }
\textcolor{blue}{
% \scriptsize
\begin{equation*}
        % \pi^* = \underset{\pi \in \Pi}{\text{argmax}} \sum_{\text{task}} f_s \left(\text{task}, F(m) | F = \arg\max_{f\in\mathcal{F}}\pi(f|\text{context}) \right) 
        \underset{\pi \in \Pi}{\text{argmax}} \sum_{\text{task}} f_s \left(\text{task}, F(m) | F = \underset{f\in\mathcal{F} }{\text{argmax}} \ \pi(f|\text{context}) \right) 
\end{equation*}
% \normalsize
}
% where $\Pi$ is the space of all possible policies, $f_s \in \mathcal{R}$ is a functional indicating the degree to which the future task (function) to be performed by the agent $\mathfrak{A}_j$ is successful.
\noindent where $\Pi$ is the space of all possible policies, $f_s(.) \in \mathcal{R}$ is a functional indicating the degree of task (function) success.
% \noindent where $\Pi$ is the space of all possible policies, $f_s(.) \in \mathcal{R}$ is a functional indicating the degree to which the task (function) to be performed is successful.
%multi-agent systems where specialized agents collaborate on complex tasks. Consider a system with agents such as a problem solver (mathematical reasoning), code executor (computational verification), and verifier (solution validation). When any agent generates a response, it invokes \method to transform the raw message into an optimal structured format before transmission to the target agent.

We define the MDP by the tuple $\langle \mathcal{S}, \mathcal{A}, \mathcal{P}, \mathcal{R}, \rangle$, where $\mathcal{S}$ is the true state space, $\mathcal{A}$ is the action space, $\mathcal{P}$ represents transition probabilities, $\mathcal{R}$ is the reward function.We thus formulate the MDP as:
(i) \textbf{State Space} $\mathcal{S}$: Task types $t \in \mathcal{T}$ representing generalizable abstraction of messages (e.g., \texttt{information\_request}, \texttt{logical\_reasoning}, \texttt{decision\_strategy})
(ii) \textbf{Action Space} $\mathcal{A}$: Message structures/representations/formats $f \in \mathcal{F}$ including structures like JSON, XML, tabular data, decision trees, etc. Note here we use the term structures to mean representations (formats) and do not limit the search space to only include ``structured'' objects (eg: JSON).
(iii) \textbf{Reward Function} $R(t,f)$: Success metric that measures goal attainment (eg: +1 for success else -1)
% \item \textbf{Reward Function} $R(t,f)$: Binary success indicator based on task completion accuracy
(iv) \textbf{Policy} $\pi(f|t)$: Probability distribution over formats/representations given task type.
% \begin{itemize}
% \item \textbf{State Space} $\mathcal{S}$: Task types $t \in \mathcal{T}$ representing communication intent (e.g., \texttt{information\_request}, \texttt{logical\_reasoning}, \texttt{decision\_strategy})
% \item \textbf{Action Space} $\mathcal{A}$: Message structures/representations/formats $f \in \mathcal{F}$ including structures like JSON, XML, tabular data, decision trees, etc. Note here we use the term structures to mean representations (formats) and do not limit the search space to only include ``structured'' objects (eg: JSON).
% \item \textbf{Reward Function} $R(t,f)$: Success metric that measures goal attainment (eg: +1 for success else -1)
% % \item \textbf{Reward Function} $R(t,f)$: Binary success indicator based on task completion accuracy
% \item \textbf{Policy} $\pi(f|t)$: Probability distribution over formats/representations given task type
% \end{itemize}

In our problem setup, we allow the tasks and the representations to evolve as the learning progresses. This is needed as in real world problems where the agents encounter new information, we may not always know the task and structures beforehand. \textcolor{blue}{In later sections (cf. fig. \ref{fig:ablation}), we show empirically that the evolving/dynamic state/action formulation benefits the performance over fixed state/action space i.e. standard MDP}. Thus we propose to evolve the MDP as the agent encounters new states (tasks) and actions (structure formats) while learning. Formally, we expand our state space $\mathcal{T}_{old}$ to $\mathcal{T}_{new} =\mathcal{T}_{old} \cup \{t_{new}\}$ on encountering new task $t_{new}$. Similarly the action space $\mathcal{F}_{old}$ is updated as $\mathcal{F}_{new} =\mathcal{F}_{old} \cup \{f_{new}\}$ on encountering new structure format $f_{new}$. As the states and actions evolve during learning, we work in the setting of an Evolving
%/Expanding 
Markov Decision Process (EMDP). One point to note is that the tasks themselves are not observed directly and our method estimates the task based on the observations. As such the assumed MDP, in theory is a partially observable one, where one might maintain some belief states from the actual environment observations, which can be intractable in real multi-agentic problem spaces. However, we observe in our experiments that the probability distribution of the states given an observation has a large mass around a state, with high accuracy (cf. $\S$ \ref{task_classification_results}). Thus we can approximately know the state and the problem, at a given time, can be approximated as an MDP to enable tractable policy learning. In the below proposition, we show the convergence of learning the proposed EMDP (proof in \textcolor{blue}{Appendix} $\S$ \ref{appendix_sec_proof}):

% \begin{proposition}[Convergence of the proposed EMDP based algorithm]\label{prop_emdp_convergence}
%     The proposed $EMDP = (S_e, A_e, P_e, R, \gamma)$, where $S_e, A_e$ are the evolving states and actions respectively, converges for the states and actions explored.
% \end{proposition}
\begin{proposition}[Convergence of the proposed EMDP based algorithm]\label{prop_emdp_convergence}
    Consider the $EMDP = (S_e, A_e, P_e, R, \gamma)$, with value function $V$. 
    % $S_e, A_e$ are the evolving states and actions respectively. 
    Let $T$ denote the update operator (including expanding $S_e, A_e$) s.t. $V_{k+1}=T(V_{k})$, then $\underset{k \rightarrow \infty}{\lim} \lVert V_k - V^* \rVert = 0$, where $V^*$ is the optimal value function for the states and actions explored till $k$ iterations.
\end{proposition}

% \md{The below 2 paragraphs need to change -- We do not need to formally define POMDP at all.}
% The underlying problem is naturally a Partially Observable Markov Decision Process (POMDP) because we maintain some estimate about the states sampled through trajectories (belief states) which might be different from the actual information about the states. We define POMDP by the tuple $\langle \mathcal{S}, \mathcal{A}, \mathcal{T}, \mathcal{R}, \Omega, \mathcal{O} \rangle$, where $\mathcal{S}$ is the true state space, $\mathcal{A}$ is the action space, $\mathcal{T}$ represents transition probabilities, $\mathcal{R}$ is the reward function, $\Omega$ is the observation space, and $\mathcal{O}$ defines observation probabilities. In our context, the true conversation state, agent internal reasoning, and complete format vocabulary remain partially observable.

% However, we approximate this POMDP as an MDP $\langle \mathcal{S}', \mathcal{A}', \mathcal{T}', \mathcal{R}' \rangle$ when the state-action space evolves slowly relative to the learning timescale. This approximation is valid when: (1) the task taxonomy stabilizes after initial exploration, (2) format discovery occurs infrequently compared to format selection decisions, and (3) the communication context can be sufficiently captured by task classification. Under these conditions, the belief state estimation complexity of POMDPs reduces to manageable state estimation in the MDP framework.

% We formulate this MDP approximation as:
% \begin{itemize}
% \item \textbf{State Space} $\mathcal{S}$: Task types $t \in \mathcal{T}$ representing communication intent (e.g., \texttt{information\_request}, \texttt{logical\_reasoning}, \texttt{decision\_strategy})
% \item \textbf{Action Space} $\mathcal{A}$: Format specifications $f \in \mathcal{F}$ including structure/representations like JSON, XML, tabular data, decision trees, etc.
% % \item \textbf{Reward Function} $R(s,a)$: Binary success indicator based on task completion accuracy
% \item \textbf{Reward Function} $R(s,a)$: Measure of success based on task completion accuracy
% \item \textbf{Policy} $\pi(a|s)$: Probability distribution over formats given task type
% \end{itemize}

% This MDP approximation enables tractable Q-learning while our cyclic expansion-convergence approach handles the residual partial observability by alternating between state-action space discovery and value function convergence.

% \subsection{Multi-Agent Integration and Learning Objective}

% Each agent invokes MetaPromptor through a \texttt{modify\_message(source, target, message)} call, creating a seamless integration where format optimization occurs transparently during normal agent communication. The system learns from complete multi-agent conversation trajectories of length $H$ ending with outcome $O \in \{0,1\}$:
% \begin{align}
% \mathbb{E}\left[\sum_{h=0}^{H} \gamma^h R_h\right] \text{ where } R_h = \alpha \cdot O - \beta
% \end{align}
% where $\alpha$ rewards correct outcomes, $\beta$ encourages efficiency, and $\gamma$ is the discount factor. Unlike traditional RL settings, our action space dynamically expands as new formats are discovered through LLM generation, requiring adaptive exploration strategies.

% \section{\method: RL-based Format Selector for Optimal Agentic Communication}
% \section{\method: Deciding Optimal Structure for Agentic Communication}
% \section{\method: Learning Optimal Message Representations for Agentic Communication}
\section{Learning Optimal Representations for Agentic Communication}

Having seen the problem setup in the previous section, this section outlines our method, \method which consists of the following components: 1) Task classification 2) Policy for optimal representation selection. The message from the source agent  $\mathfrak{A}_i$ is categorized into the appropriate task. This is done using an LLM, since LLMs have exhibited capabilities for few/zero shot text classification tasks \cite{brown2020languagemodelsfewshotlearners}. Once the task is decided, the learnt policy then selects the optimal structure using which the message is formatted for communication with destination agent $\mathfrak{A}_j$. The policy is learnt using an adapted policy learning framework as we shall see in the next sections. The system overview for \method is in figure \ref{fig:system_overview}.
% The system operates through two main LLM-based components: task classification and format selection, followed by Q-learning updates from trajectory outcomes.

\begin{figure*}[t!]
    \centering
    \includegraphics[width=0.95\linewidth]{AAAI/figures/system_overviewV2.png}
    % \caption{Framework Overview: Left figure shows the overview of our method which consists of two components namely the task categorization and the optimal representation policy learning module. The optimal representation module learns the policy for the chosen task using historical data or online from the reward on achieving the end goal. Right figure compares the outcomes of multi-agent communication on a sample problem, using base and our method. We see that the system without our method chooses natural language and arrives at the wrong answer whereas our method selects code as the message representation to deduce the correct answer.}
    \caption{Left figure shows the overview of our method which consists of two components namely the task categorization and the policy learning module (in blue). The policy module learns the policy for the chosen task using historical data or online from the reward on achieving the end goal. Right figure compares the outcomes on a sample problem, using base and our method. The vanilla system chooses natural language and arrives at the wrong answer whereas our method selects code as the message representation to deduce the correct answer.}
    \label{fig:system_overview}
\end{figure*}

% \subsection{State and Action Selection via LLM Prompting}
% \subsection{Task categorization}

\noindent \textbf{\textcolor{blue}{1) Initialization:}}
\textcolor{blue}{In order to initialize the task and representations, we sample and process the data (if available) to build a knowledge base (KB), by prompting the LLM. In case of a pure online setting, we bootstrap the KB from prior knowledge of data from similar domain. The novel tasks and representations are then updated as the learning progresses. Note, we do not aim for OOD generalization but adaptation to the concerned task/domain.
}

\noindent \textbf{2) Task categorization:}
Tasks serve as an abstraction to help generalize to novel messages. Prior to selecting the optimal representation using the policy, the system needs to classify the incoming message into one of the known tasks or a new one. For this categorization, we use an LLM. We warm start the task taxonomy with a few known categories and adopt a chain of thought style prompting, where the LLM is prompted to first review the message and decide whether the message fits into one of the existing tasks or a new task category needs to be added to the set, followed by the categorization of the message into the task. Following are the inputs to the LLM:
(i) Source agent name and description (e.g., ``problem\_solver: mathematical reasoning agent'')
(ii) Target agent name and description (e.g., ``code\_executor: computational verification agent'') 
(iii) Raw message content requiring formatting
(iv) Current task taxonomy with descriptions.

% \begin{itemize}
% \item Source agent name and description (e.g., ``problem\_solver: mathematical reasoning agent'')
% \item Target agent name and description (e.g., ``code\_executor: computational verification agent'') 
% \item Raw message content requiring formatting
% \item Current task taxonomy with descriptions
% \end{itemize}
Thus the module considers the message as well as the agent context in order to decide the task category. 
\textcolor{blue}{If the existing task categories do not fit the current message context, the module (LLM) creates a new task category.}
The prompt for selection of task is provided in the \textcolor{blue}{Appendix $\S$ \ref{prompts}}. Once the task is decided, we then provide the task to the policy which learns and decides the optimal representation for that task and agent pair.

% \paragraph{Task Classification (State Selection)}
% When \texttt{get\_format()} is invoked, \method first identifies the communication task type through an LLM call. The prompt includes:
% \begin{itemize}
% \item Source agent name and description (e.g., "problem\_solver: mathematical reasoning agent")
% \item Target agent name and description (e.g., "code\_executor: computational verification agent") 
% \item Raw message content requiring formatting
% \item Current task taxonomy with descriptions
% \end{itemize}

% Depending on the learning phase within an epoch, the system either: (1) \textbf{Expansion Phase}: allows the LLM to generate novel task types, or (2) \textbf{Convergence Phase}: constrains selection to existing task categories to ensure Q-value convergence. The exact prompt for task-identification is given in appendix.

% \paragraph{Format Selection (Action Selection)}
% \paragraph{Optimal Structure Representation Discovery}
% \subsection{Deciding Message Representation}
\noindent \textcolor{blue}{
\textbf{3) Policy for Representation Selection:}
We have seen the formulation of the problem-setting as expanding MDP. This module learns a policy ($\pi$) to sample from the action (structure) distribution given the state (task). In order to learn the optimal policy $\pi^*$, we design a policy learning method for our problem setting. The algorithm is an off policy method to allow for exploration using a behavior policy $\mu(a|s)$.
% to sample trajectories (at inference we sample greedily from the learnt optimal policy).
At inference we sample greedily from the learnt optimal policy.
Exact methods would require an exponential complexity to compute the system dynamics (cf. $\S$ \ref{intro}). Thus we approach the solution using trajectories as the agents explore the environment and employ Monte Carlo (MC) based methods \cite{sutton2018reinforcement}.
}
% \noindent \underline{(I) \emph{Deciding Message Representations:}}
% At inference, 
% the optimal representation module takes a task (given by the task categorization module) and decides the structure in which to format the message for downstream communication; formulating the problem-setting as expanding MDP. Internally, the module learns a policy ($\pi$) to sample from the action (structure) distribution given the state (task). In order to learn the optimal policy $\pi^*$, we design a policy learning method for our problem setting. During inference, we desire to adopt a simple $\epsilon$-greedy policy that selects the best action $\pi(a|s)$. Thus the algorithm is an off policy method to allow for exploration using a behavior policy $\mu(a|s)$ to sample trajectories. As we have seen exact methods would require an exponential complexity to compute the system dynamics, we rather approach the solution using empirical trajectories as the agents explore the environment. As such, we employ Monte Carlo (MC) based methods \cite{sutton2018reinforcement} for the learning algorithm.

% \subsection{Discovery of Optimal Representations}
% \noindent \textbf{3) Discovery of Tasks and Representations:}
% \noindent \textbf{\textcolor{blue}{4) Discovery of Representations:}}
\noindent \underline{\textcolor{blue}{(3.1) \emph{Discovery of Representations:}}}
In our method, the behavior policy explores the tasks and representations followed by policy optimization. We first describe the behavior policy before we see how we use it in our algorithm to find the optimal policy ($\pi^*(a|s)$). The behavior policy $\mu(a|s)$ consists of sampling using three strategies namely: 1) Sampling from Q values (Q) 2) LLMs as a sampler (L) and 3) Diverse Exploration (D).
The motivation for using these policies comes from the seminal work by Freud \cite{freud1989ego}, which suggests that human behavior is influenced by unconscious memories/patterns and rational thoughts. Inspired by human behavior, the behavior policy samples actions from its ``unconscious'' learnt patterns (Q-values), ``conscious'' rational module (LLM) and random exploration (diversity policy).
% In order to give differential preference to each of the samplers, as the learning progresses, we use adaptive weights $\mathbf{w} = [w_Q, w_L, w_D]$ for each strategy respectively. Below we elaborate on each of the strategy.
% In order to give differential preference to each of the samplers, we use weights $\mathbf{w} = [w_Q, w_L, w_D]$ for each strategy respectively. 
In order to give differential preference to each of the samplers, we use strategy-wise weights $\mathbf{w} = [w_Q, w_L, w_D]$. 
% Below we elaborate on each of the strategy.

\noindent \textbf{Sampling from Q values (Q)}: 
This policy samples structures from the learnt Q values, akin to a scoreboard, in an $\epsilon-$greedy manner. Specifically, the probability of the action (message structure) is computed as:
$\pi_Q(a|s) = \frac{\exp(Q(s,a)/\tau)}{\sum_{a'} \exp(Q(s,a')/\tau)}$ if $\text{random()}>\epsilon$ else $\frac{1}{\lvert \mathcal{F} \lvert}$.
where $\tau$ is the sampling temperature. Thus this policy encourages message structures that are known to work well for a task.
% This policy simply exploits the learnt Q values to perform sampling. Specifically, a softmax is taken over the Q values as follows
% $\pi_Q(a|s) = \frac{\exp(Q(s,a)/\tau)}{\sum_{a'} \exp(Q(s,a')/\tau)}$
% where $\tau$ is the sampling temperature.

\noindent \textbf{LLMs as a sampler (L)}: This policy selects the action (message structure) by prompting an LLM with comprehensive context:
1) Source and target agent context descriptions
2) Message content and categorized task type
3) Historical pass rates (\% success) for each structure previously used with this task
4) Currently overused formats to encourage diversity
% 5) High-level optimization goals: ``enhance clarity, increase brevity, provide structured formats that represent logical relationships, reduce token usage ...''.
5) High-level optimization goals (e.g. enhance clarity, efficiency etc.).
% \item High-level optimization goals: "enhance clarity, increase brevity, provide structured formats that represent logical relationships, reduce token usage, enable systematic analysis"
% \begin{itemize}
% \item Source and target agent context descriptions
% % \item Source and target agent descriptions for context-aware formatting
% \item Message content and categorized task type
% \item Historical pass rates (\% success) for each structure previously used with this task
% \item Currently overused formats to encourage diversity
% \item High-level optimization goals: ``enhance clarity, increase brevity, provide structured formats that represent logical relationships, reduce token usage ...''
% % \item High-level optimization goals: "enhance clarity, increase brevity, provide structured formats that represent logical relationships, reduce token usage, enable systematic analysis"
% \end{itemize}
% This is analogous to hard-mining,
% in Dense-Passage-Retriever (DPR) \cite{dpr}, 
In doing so, we look for representations that are specifically designed with the current query in mind. 
\textcolor{blue}{This phase also allows for addition of new message structures as needed and decided by the LLM.
The prompt for structure selection via LLM is given in $\S \ref{prompts}$.}

% \textcolor{blue}{
\noindent \textbf{Diverse Exploration (D)}: 
% In addition to exploiting the Q-values learnt or relying on the inductive biases of the LLM to sample message representations for a task, we would also like the method to explore diverse structures not yet encountered in previous trajectories. 
Here, we explore diverse formats that did not appear for concerned task but appeared for other tasks. To give a fair chance to the under-sampled formats, we sample inversely by frequency.
% explore diverse structures not yet encountered in previous trajectories. 
Thus we design this policy which samples inversely to usage frequency for exploration across tasks from the distribution $p_i = \frac{(1- n_i/N)}{(N-1)}$, 
% where $p_i$ is the sampling probability for task $i$, $n_i$ is the count of occurrence of each task and $N$ is the total steps. 
where $p_i, n_i$ are the sampling probability and count for task $i$ and $N$ is the total steps. 
% This is analogous to soft-mining in DPR.

% The strategy weights are updated using softmax over performance: $w_i^{(t+1)} = \frac{\exp(\bar{R}_i^{(t)})}{\sum_{j} \exp(\bar{R}_j^{(t)})}$, where $\bar{R}_i^{t}$ is the average reward received till time $t$ using policy $i$.
% }
% \subsection{Off-Policy Monte Carlo Learning with Environment Expansion}
% \paragraph{Off-Policy Monte Carlo Learning with Environment Expansion}
% \noindent \textbf{5) Learning the Optimal Policy:}
\noindent \textcolor{blue}{\underline{(3.2) \emph{Learning the Optimal Policy:}}}
Having described the behavior policy, we now explain our algorithm in this section. As seen, we adopt an off policy method in which the states and actions are not fixed. 
We adopt a Monte Carlo (MC) based learning since we obtain feedback at the end of the trajectory. MC methods are known to be unbiased, stable and avoid spurious noisy rewards.
Specifically, the method can be described as an off-policy MC learning with environment expansion.
The target policy $\pi(a|s)$ (for inference) greedily selects structures. 
Formally,
% \begin{align}\label{target_policy_eq}
% \begin{split}
%   \pi(a|s) &= 1 \quad \text{if} \quad a = \arg\max_{a'} Q(s,a')\\ 
%            &= 0 \quad \text{otherwise}
% \end{split}
% \end{align}
\begin{align}\label{target_policy_eq}
\pi(a|s) = \begin{cases}
  1 & \text{if } a = \arg\max_{a'} Q(s,a') \\ 
  0 & \text{otherwise}
\end{cases}
\end{align}
% $\pi(a|s) = 1$ if $a = \arg\max_{a'} Q(s,a')$, else $0$.
The behavior policy combines the exploration strategies as follows:
\textcolor{blue}{
% {\small
\begin{equation*}\label{behavior_policy_eq}
    \mu(a|s) = w_Q \pi_Q(a|s) + w_L \pi_L(a|s) + w_D \pi_D(a|s)
\end{equation*}
% }
}\vspace{-0.5cm}
% $\mu(a|s) = w_S \pi_S(a|s) + w_L \pi_L(a|s) + w_R \pi_R(a|s)$.

The Q-values are updated using importance sampling, on every visit as per the below expression:
\vspace{-0.5cm}\textcolor{blue}{
% {\small
\begin{align*}\label{q_update}
Q_{t+1}(s,a) &\leftarrow Q_t(s,a) + \frac{W_{t+1}}{C_t(s,a)}(G_t - Q_t(s,a)) \\
G_{t+1} &= \gamma G_t + R_{t+1}, \quad W_{t+1} = \frac{W_t}{\mu(a|s)}
\end{align*}
% }
}
where $Q \in \mathcal{R^{\lvert T \lvert \times \lvert F \lvert}}$ is the Q table, $W$ is the importance sampling ratio, $R_{t+1}$ is the reward at time $t$ for taking action $a \in \mathcal{R^F}$ (message structure) in state $s \in \mathcal{R^T}$ (task), $C_t(s,a)$ is a normalizing factor (generally count of the times $s,a$ is visited).

In the above, the tasks (state) and structures (action) could evolve as the learning progresses.
Since new state-action pairs are discovered during learning, we employ the below two phases in learning:

\textbf{Expansion Phase}: Allow discovery of new $(s,a)$ pairs up to threshold $\alpha$ fraction of interactions. In this phase the LLM generates new tasks categories and structures as seen previously.

\textbf{Convergence Phase}: 
% Once new tasks and structures appear infrequently, we fix the state-action space and allow for Q-value convergence. 
Once new tasks and formats become rare, we fix the state-action space and allow for Q-value convergence. 

% For deployment, the system samples greedily from learned Q-values, using the target policy $\pi(a|s)$ for optimal performance.

% The off-policy Monte-Carlo Algorithm is adapted from \cite{sutton2018reinforcement} (and given in Appendix) and is not a novel contribution of the paper. The novel contribution is the design of \method and the design of behaviour-policy $\mu$

The system learns from the complete multi-agent conversation trajectories of length $L_T$ ending with outcome $R \in \{0,1\}$. Algorithm \ref{appendix:offpolicy_mc} in appendix summarizes the method. The complexity of the method is $\mathcal{O}(\mathcal{(\lvert T \rvert + \lvert F \rvert) \lvert D \rvert}L_T)$ (cf. $\S$ \ref{computational_complexity}), which is linear in number of tasks ($\mathcal{T}$), formats ($\mathcal{F}$), dataset size ($\mathcal{D}$) and trajectory length ($L_T$).

% \begin{algorithm}[h]
% % \caption{Off-Policy Every-Visit MC Control Algorithm Using Weighted Importance Sampling}
% % \caption{Off-Policy Every Visit Monte Carlo Learning using Importance Sampling with Environment Expansion}
% \caption{\method Algorithm}
% \label{alg:offpolicy_mc}
% \begin{algorithmic}[1]
% \STATE \textbf{Initialize}, for all $s \in \mathcal{T}$, $a \in \mathcal{F}(s)$:
% \STATE \quad $Q(s,a) \leftarrow$ arbitrary
% \STATE \quad $C(s,a) \leftarrow 0$
% \STATE \quad $\mu(a|s) \leftarrow w_Q \pi_Q(a|s) + w_L \pi_L(a|s) + w_D \pi_D(a|s)$ (behavior policy)
% \STATE \quad $\pi(s) \leftarrow$ a deterministic policy that is greedy wrt $Q$
% % \STATE
% \REPEAT
%     % \STATE //Note in the below the policy $\pi_L$ (component of $\mu$) will add new tasks and structures as needed and expand the state action space
%     \STATE msg $\leftarrow$ Sample from the set of queries/problems
%     \STATE Generate an episode using behavior policy $\mu(a|s)$:
%         \REPEAT 
%             \STATE $\mathcal{T}_{t} \leftarrow \text{TaskCategorization(msg, context)}$
%             \STATE $\mathcal{F}_t \leftarrow \mu(T_{t}, context)$
%             \STATE prompt $\leftarrow$ ApplyStructure(msg, $\mathcal{F}_t$)
%             \STATE msg, $R_{t+1} \leftarrow \mathfrak{A}_t$(prompt)
%             % \STATE // New tasks and structures are added during expansion phase
%             \IF{Expansion Phase} 
%                 \STATE $\mathcal{T} \leftarrow \mathcal{T} \cup \{\mathcal{T}_{t}\}$, $\mathcal{F} \leftarrow \mathcal{F} \cup \{F_t\}$
%             \ENDIF
%             \STATE $t \leftarrow t+1$
%         \UNTIL{Episode end}
%     % \STATE \quad $\mathcal{T}_0, \mathcal{F}_0, R_1, \mathcal{T}_1, \mathcal{F}_1, R_2, \ldots, \mathcal{T}_{T-1}, \mathcal{F}_{T-1}, R_T, S_T$
%     \STATE $G \leftarrow 0$
%     \STATE $W \leftarrow 1$
%     \FOR{$t = T-1, T-2, \ldots$ down to $0$}
%         \STATE $G \leftarrow \gamma G + R_{t+1}$
%         \STATE $C(\mathcal{T}_t, \mathcal{F}_t) \leftarrow C(\mathcal{T}_t, \mathcal{F}_t) + W$
%         \STATE $Q(\mathcal{T}_t, \mathcal{F}_t) \leftarrow Q(\mathcal{T}_t, \mathcal{F}_t) + \frac{W}{C(\mathcal{T}_t, \mathcal{F}_t)}[G - Q(\mathcal{T}_t, \mathcal{F}_t)]$
%         \STATE $\pi(\mathcal{F}_t|\mathcal{T}_t) \leftarrow 1 \quad \text{if} \quad  \mathcal{F}_t=\arg\max_a Q(\mathcal{T}_t, a)$ else 0 (with ties broken arbitrarily)
%         \STATE $W \leftarrow W \cdot \frac{\pi(\mathcal{F}_t|\mathcal{T}_t)}{\mu(\mathcal{F}_t|\mathcal{T}_t)}$
%         \IF{$W = 0$}
%             \STATE \textbf{Exit For Loop}
%         \ENDIF
%     \ENDFOR
% \UNTIL{Convergence or max iterations}
% \end{algorithmic}
% \end{algorithm}
% \begin{align}
% \mathbb{E}\left[\sum_{h=0}^{H} \gamma^h R_h\right] \text{ where } R_h = \alpha \cdot O - \beta
% \end{align}
% where $\alpha$ rewards correct outcomes, $\beta$ encourages efficiency, and $\gamma$ is the discount factor. 
% Unlike traditional RL settings, our action space dynamically expands as new formats are discovered through LLM generation, requiring adaptive exploration strategies.


\noindent \textcolor{blue}{\underline{(3.3) \emph{Inferring Representations:}}}
\textcolor{blue}{At inference, 
the policy module takes a task (given by the task categorization module) and decides the structure in which to format the message for downstream communication. We adopt a greedy policy (eq. \ref{target_policy_eq}) that selects the best action $\pi^*(a|s)$.}

% % \subsection{Integration in Multi-Agent Systems}
% \noindent \textbf{4) Integration in Multi-Agent Systems: }
% During inference, the method categorizes the message into one of the 
% % frozen 
% task taxonomies and samples the message structure greedily from the learned Q-values as in equation \ref{target_policy_eq}, using the target policy $\pi(a|s)$ for optimal performance.
% Each agent invokes \method through a \texttt{modify\_message(source, target, message)} call, 
% for optimizing message representation.
% % creating a seamless integration where format optimization occurs transparently during agent communication. 
% % The system learns from complete multi-agent conversation trajectories of length $H$ ending with outcome $O \in \{0,1\}$:
% % \begin{align}
% % \mathbb{E}\left[\sum_{h=0}^{H} \gamma^h R_h\right] \text{ where } R_h = \alpha \cdot O - \beta
% % \end{align}
% % where $\alpha$ rewards correct outcomes, $\beta$ encourages efficiency, and $\gamma$ is the discount factor. Unlike traditional RL settings, our action space dynamically expands as new formats are discovered through LLM generation, requiring adaptive exploration strategies.



\section{Experiment Setup}
We evaluate our method \method, following the approach of the main baseline \cite{chen-etal-2024-beyond-natural}. 
Specifically, we seek to answer the following research questions:
(i) \textbf{RQ1}: Does \method enhance the performance over static and fixed structures in multi-agent settings?
(ii) \textbf{RQ2}: Do we need to optimize for representations or are there fixed structures that work across datasets?
(iii) \textbf{RQ3}: How is the efficiency of the system affected?
    % \item \textbf{RQ4}: What is the effectiveness of each component of \method?
(iv) \textbf{RQ4}: What is the effectiveness of each component?
(v) \textbf{RQ5}: What is the distribution of the dynamically evolving tasks and representations?
% \begin{itemize}
%     % \item \textbf{RQ1}: Does \method enhance the performance over static and fixed structures in multi-agent communication settings?
%     \item \textbf{RQ1}: Does \method enhance the performance over static and fixed structures in multi-agent settings?
%     \item \textbf{RQ2}: Do we need to optimize for representations or are there fixed structures that work work across datasets?
%     % \item \textbf{RQ3}: How does \method affect the efficiency of the system?
%     \item \textbf{RQ3}: How is the efficiency of the system affected?
%     % \item \textbf{RQ4}: What is the effectiveness of each component of \method?
%     \item \textbf{RQ4}: What is the effectiveness of each component?
%     \item \textbf{RQ5}: What is the distribution of the dynamically evolving tasks and structures?
% \end{itemize}
% we benchmark our method \method using the setup as detailed below:
We detail the experiment setup below

\subsection{Datasets:}
\noindent \textbf{Single LLM Reasoning:} In the motivation study (cf. fig. \ref{fig:motivation}), we employ the datasets of Logic Grid, Coin Flip and AQuA adopting the same experiment setting as \cite{chen-etal-2024-beyond-natural}. 
% (cf. appendix \ref{appendix_single_agent}).

\textbf{Multi-Agent:}
(i) \textbf{Collaborative Problem Solving:}
We use the GSM+ (GSM) benchmark~\cite{gsmplus} to evaluate mathematical problem solving, as the adversarial examples in the dataset could benefit from a multi-agent setup and various message structures. 
% We design a multi-agent scenario with three agents---\textbf{Planner}, \textbf{Verifier}, and \textbf{Coder}---to address complex reasoning tasks. 
% We also evaluate on additional reasoning datasets (to be specified). 
% Correctness is measured using exact string match.
Correctness is measured by matching answers upto precision.
% Correctness is measured using standard metrics.
% for each benchmark. 
% Following \cite{cemri2025multi-agent}, we analyze performance along standard failure modes for LLM-based multi-agent systems, using an LLM as judge. 
% We track interaction metrics such as the average response length.
% number of conversational turns and average response length.
(ii) \textbf{Information Exchange}
In order to evaluate the ability to find optimal structures, we experiment with information exchange using multi-hop QA using the NarrativeQA (NQA)~\cite{narrativeqa}, WikiHopQA (WQA)~\cite{wikihop}, and HotpotQA (HQA)~\cite{hotpotqa} datasets. Correctness is evaluated as follows: (i) when answer choices are provided, we use exact match, (ii) for open-ended responses, we employ an NLI-based metric where we use a RoBERTa model fine-tuned for NLI \cite{roberta_nli}, treating the question and model-generated answer as the ``premise,'' and the annotation as the ``hypothesis,'' scoring correct if entailment score > 0.7 (set empirically). 
% (threshold set empirically). 
% We thus capture we are interested in semantic match as compared to LCS-based surface-level match.
% We do not follow the approach of \cite{chen-etal-2024-beyond-natural} to utilise RougeL as metric since we are interested in semantic match as compared to LCS-based surface-level match. 
% As above, we monitor common failure modes~\cite{cemri2025multi-agent} and interaction metrics.

% \textbf{3. FLASH Multi-Agentic System for Triaging}
% FLASH is a multi-agent triage framework designed to enhance the speed, accuracy, and adaptability of incident prioritization in dynamic environments. Leveraging a collection of specialized and collaborating agents, FLASH facilitates automated triage by integrating semantic analysis, dynamic knowledge sharing, and negotiation mechanisms across agents. This architecture improves response times and ensures consistent decision-making, particularly in high-stakes scenarios such as emergency services or cloud incident management. FLASH has demonstrated improved triage accuracy and operational efficiency over traditional manual or rule-based systems~\cite{yu2025flash}.

\subsection{Baselines:}
For all the experiments, we choose the vanilla multi-agent system (Vanilla MAS) and Autoform (prompts LLM for best structure from predefined ones) as the baselines. 
Please note that the Vanilla MAS already uses prompting techniques like CoT \cite{cot}, ReAct \cite{react_iclr} etc. and our method is orthogonal to these. 
Since these methods involve directly prompting the LLMs and there is no learning involved, we directly evaluate on the test-set. For \method, we first learn the policy by running trajectories over a small sample of data points ($\sim 500-2000$). The testing set is obtained keeping the similar setting as baseline Autoform ~\cite{chen-etal-2024-beyond-natural}. In order to gauge the benefits of dynamically learning the message structures from data/environment, we experiment with LLMs such as the multimodal GPT4o \cite{openai2024gpt4o} and large reasoning models such as O3 \cite{openai2025competitiveprogramminglargereasoning}. To further understand the effect of learning structures on smaller models, we also experiment with some small language models such as Phi4-mini \cite{phi4technicalreport}, llama3-8B \cite{llama3} and Qwen2.5-Math-7B \cite{qwen2025qwen25technicalreport}.

\subsection{Implementation details}
% \textcolor{blue}{
We design multi-agent scenarios with three agents - Planner, Verifier, and Coder - for GSM+ dataset and five agents (each with different context/information) for the multi-hop setting, using the AutoGen framework \cite{wu2023autogen}. 
For the policy learning, we initialize the Q values for newly discovered structures and sampling Temperature as 0. We allow the trajectories to roll out for a maximum of 25 steps (as we empirically find the goal is achieved by then). For the behavior policy, we set the weights as $w_Q=0.5, w_L=0.25, w_D=0.25$. We set the $\epsilon=0.1$.
% equal values and then depending on the rewards we adapt them dynamically at run time. 
The expansion phase is allowed for around 20\% of the epoch after which we allow the Q-values to converge. The trajectory is terminated when any agent thinks the goal is achieved or the maximum number of iterations (25) have been reached. If the goal/answer matches the ground truth a reward of +1 is provided else a penalty of -1 is incurred, for every step.
All experiments are conducted on a linux machine with a single A100 GPU. For SLM-inference, we use the vLLM library \cite{vllm} and official open source models as on huggingface \cite{huggingface}. 
Each experiment is run three times and we report the average over runs for brevity (variance $\sim \pm 2\%$). 
% }


\section{Results}

% As we have already seen in the motivation study (cf. fig. \ref{fig:single_agent}), learning optimal representations to prompt the LLM outperforms the vanilla counterpart using natural language message representation. In this section we extend the experiments and present results on the multi-agentic communication setup to demonstrate the efficacy of learning optimal message representations in this setting.
% 1. Incorporating real-world datasets such as WebArena, GAIA etc.
% 2. Incorporating situations with partial-observability
% 3. Complex Debates (mejority voting, reward-guided tree search, latent space reasoning via step distillation)
\begin{table}[t!]
  \caption{\textcolor{blue}{\method outperforms the vanilla MAS using NL and Autoform (fixed formats by the LLM).}} 
    \label{tab:main_results}
    \centering
    % \begin{tabular}{c{0.2cm} c{0.2cm} c{0.2cm} c{0.2cm} c{0.2cm} c{0.2cm} c{0.2cm} c{0.2cm} c{0.2cm}}
    % \resizebox{\columnwidth}{!}{
    \resizebox{0.5\textwidth}{!}{
    \begin{tabular}{c c c c c c c}
      % \textbf{Base LLM} & \textbf{Method} & \textbf{GSM+} & \textbf{WikiHop} & \textbf{HotpotQA} & \textbf{NarrativeQA} & \textbf{Average} \\
      \textbf{LLM} & \textbf{Method} & \textbf{GSM} & \textbf{WQA} & \textbf{HQA} & \textbf{NQA} & \textbf{Avg} \\
        \hline \hline
        % GPT-4o & Vanilla MAS & 86.5 & 10.2 & 60.7 & 47.1 & 51.125 \\
        % & Autoform & 87.5 & 11.5 & 62.5 & 48.6 & 52.525 \\
        % & \method (Ours) & 89.5 & 13.9 & 64.5 & 51.1 & 54.75 \\
        % \hline
        % o3 & Vanilla MAS & 93.5 & 15.5 & 65.2 & 53.6 & 56.95 \\
        % & Autoform & 93.9 & 16.3 & 66.1 & 55.8 & 58.02 \\
        % & \method (Ours) & 94.7 & 17.2 & 67.3 & 56.7 & 58.98 \\
        % \hline
        % Phi4-mini & Vanilla MAS & 49.2& 5.5& 29.6& 22.9& 26.8 \\
        % & Autoform & 50.7& 6.2& 31.4& 23.8& 28.02 \\
        % & \method (Ours) & 53.5& 8.3& 34.2& 25.9& 30.48 \\
        % \hline   
        GPT-4o & Vanilla & 86.5 & 19.5 & 65.1 & 51.3 & 55.6 \\
        & Autoform & 87.5 & 18.2 & 48.5 & 35.9 & 47.5 \\
        & \method & \textbf{89.5} & \textbf{21.4} & \textbf{69.1} & \textbf{57.2} & \textbf{59.3} \\
        \hline
        o3 & Vanilla & 93.5 & 20.5 & 80.2 & 70.5 & 66.2 \\
        & Autoform & 92.1 & 21.9 & 81.7 & 59.8 & 63.9 \\
        & \method & \textbf{94.7} & \textbf{22.8} & \textbf{88.9} & \textbf{71.9} & \textbf{69.6} \\
        % \hline
        % Phi4- & Vanilla & 49.2 & 8.3 & 24.1 & 21.5 & 25.8 \\
        % mini& Autoform & 50.7 & 12.2 & 31.5 & 12.2 & 26.6 \\
        % & \method & 53.5 & 15.7 & 33.5 & 25.5 & 32.1 \\
        \hline
        
        % Llama-3-8B & Vanilla MAS & 93.5 & 15.5 & 65.2 & 53.6 & 56.95 \\
        % & Autoform & 93.9 & 16.3 & 66.1 & 55.8 & 58.02 \\
        % & \method (Ours) & 94.7 & 17.2 & 67.3 & 56.7 & 58.98 \\
        % \hline
        % Mistral-7B & Vanilla MAS & 93.5 & 15.5 & 65.2 & 53.6 & 56.95 \\
        % & Autoform & 93.9 & 16.3 & 66.1 & 55.8 & 58.02 \\
        % & \method (Ours) & 94.7 & 17.2 & 67.3 & 56.7 & 58.98 \\
        \hline \hline 
        Qwen2.5- & Vanilla & 79.1 & 11.8 & 29.3 & 26.2 & 36.6 \\
        7B (Math) & \method & \textbf{82.5} & \textbf{17.2} & \textbf{37.4} & \textbf{29.7} & \textbf{41.7} \\
        \hline 
        Phi4- & Vanilla & 64.5 & 6.3 & 24.5 & 22.4 & 29.4 \\
        mini & \method & \textbf{70.5} & \textbf{11.6} & \textbf{31.8} & \textbf{23.9} & \textbf{34.4} \\
        \hline 
        Llama3- & Vanilla & 77.9 & 9.7 & 28.3 & 25.1 & 35.2 \\
        8B & \method & \textbf{83.5} & \textbf{12.8} & \textbf{33.6} & \textbf{27.2} & \textbf{39.3} \\
        \bottomrule
        % \hline
        \textcolor{blue}{Efficiency} & \textcolor{blue}{$\Delta$ Tokens (\%)} & \textcolor{blue}{-8.7} & \textcolor{blue}{-5.4} & \textcolor{blue}{-18.8} & \textcolor{blue}{19.3} & \textcolor{blue}{-3.4} \\
        % \bottomrule
        \hline
        % \hline
        % \hline
        \bottomrule
    \end{tabular}
    }
\end{table}


\subsection{Main Results (RQ1)}

In this section, we present the results of our method w.r.t baselines on the benchmark datasets. Table \ref{tab:main_results} shows the accuracy results on GSMPlus and the MultiHop datasets, against the respective LLMs.
% % \begin{wraptable}{L}{3.3cm}
% % \caption{\small{Results on SLM}}\label{tab:slm_results}
% % \vspace{-0.5cm}
% % \resizebox{0.22\textwidth}{!}{
% % \begin{tabular}{cccc}\\
% %         % \toprule 
% %       \hline
% %       \hline
% %       \\
% %       \textbf{SLM} & \textbf{Method} & \textbf{WQA} & \textbf{HQA} \\\midrule
% %         \hline \hline 
% %         Qwen2.5- & Vanilla & 6.9 & 28.7 \\
% %          7B (Math) & Ours & \textbf{7.2} & \textbf{32.5} \\ 
% %          % \midrule
% %          \hline
% %         Phi4- & Vanilla & 5.3 & 24.4 \\
% %         mini & Ours & \textbf{12.1} & \textbf{31.3} \\ 
% %          % \midrule
% %          \hline
% %         Llama3- & Vanilla & 1.7 & 23.4 \\
% %         8B & Ours & \textbf{5.3} & \textbf{25.8} \\ 
% %         % \bottomrule
% %         \hline
% %         \hline
% % \end{tabular}
% % }
% % \end{wraptable} 
% \begin{wraptable}{L}{4.2cm}
% \caption{\small{Results using SLMs}}\label{tab:slm_results}
% \vspace{-0.5cm}
% \resizebox{0.28\textwidth}{!}{
% \begin{tabular}{cccccc}\\
%         % \toprule 
%       \hline
%       \hline
%       \\
%       \textbf{SLM} & \textbf{Method} & \textbf{GSM} & \textbf{WQA} & \textbf{HQA} & \textbf{NQA}\\\midrule
%         \hline \hline 
%         Qwen2.5- & Vanilla & 79.1 & 11.8 & 29.3 & 26.2 \\
%         7B (Math) & Ours & \textbf{82.5} & \textbf{17.2} & \textbf{37.4} & \textbf{29.7} \\
%         \hline 
%         Phi4- & Vanilla & 64.5 & 6.3 & 24.5 & 22.4 \\
%         mini & Ours & \textbf{70.5} & \textbf{11.6} & \textbf{31.8} & \textbf{23.9} \\
%         \hline 
%         Llama3- & Vanilla & 77.9 & 9.7 & 28.3 & 25.1 \\
%         8B & Ours & \textbf{83.5} & \textbf{12.8} & \textbf{33.6} & \textbf{27.2} \\
%         % \bottomrule
%         \hline
%         \hline
% \end{tabular}
% }
% \end{wraptable} 
We observe that our method outperforms the baselines and the enhancements are statistically significant (using Wilcoxon signed-rank test, p-value $< 0.05$). This indicates that learning the optimal structures dynamically enhances performance in multi-agent setting, answering RQ1. 
\textcolor{blue}{Where Roberta based NLI is used for evaluation, we perform a human evaluation on a sample ($\sim 100$) set and find a high Cohen's Kappa \cite{Cohen07} (>90\%).}
We also find our method (vanilla), on average, needs around 3.5 (3.5), 2.5 (3) hops using GPT4o, o3 respectively. 
% We further perform experiments on SLMs to show performance enhancement w.r.t. the vanilla multi-agent system. From the results,
From the SLM results,
% in table \ref{tab:slm_results}, 
 we can observe that our method demonstrates consistent improvements compared to the vanilla method, confirming the benefit of communicating using optimal message structures even in small models. \textcolor{blue}{We further report the percent difference of tokens consumed by our method (w.r.t. vanilla system) across methods and datasets and find that the communication cost is on par or even reduced in some cases.}

\noindent \textbf{Task Classification:}\label{task_classification_results} One of the important modules in our method is the task classification module, where the LLM samples a state from its observation (input message). Here we quantify the accuracy of the task classification, by a double-blind human review. We employ two human experts
% \footnote{Authors perform review of the tasks classified by the LLM mixed with synthetic tasks to prevent bias} 
 to label the tasks ($\sim 100$ per dataset) as classified by the module. The accuracy (across datasets) was 99.5\% and the Cohen's Kappa \cite{Cohen07} was near 100\% showing complete inter-rater agreement.  

% \begin{table}[h!]
%   \caption{Results on SLMs.} 
%     \label{tab:slm_results}
%     \centering
%     % \begin{tabular}{c{0.2cm} c{0.2cm} c{0.2cm} c{0.2cm} c{0.2cm} c{0.2cm} c{0.2cm} c{0.2cm} c{0.2cm}}
%     % \resizebox{\columnwidth}{!}{
%     \resizebox{0.5\textwidth}{!}{
%     \begin{tabular}{c c c c}
%       \textbf{SLM} & \textbf{Method} & \textbf{WQA} & \textbf{HQA} \\
%         \hline \hline 
%         Qwen2.5-7B & Vanilla & 6.9 & 28.7 \\
%          & Ours & 7.2 & 32.5 \\ 
%         Phi4-mini & Vanilla & 5.3 & 24.4 \\
%          & Ours & 12.1 & 31.3 \\
%         Llama3-8B & Vanilla & 1.7 & 23.4 \\
%          & Ours & 5.3 & 25.8 \\
%         \hline
        
%         \hline
%     \end{tabular}
%     }
% \end{table}

% \begin{wraptable}{r}{5.5cm}

% \begin{wraptable}{L}{5.5cm}
% \caption{Results on SLMs.}\label{tab:slm_results}
% \resizebox{0.3\textwidth}{!}{
% \begin{tabular}{cccc}\\\toprule  
%       \textbf{SLM} & \textbf{Method} & \textbf{WQA} & \textbf{HQA} \\\midrule
%         \hline \hline 
%         Qwen2.5-7B & Vanilla & 6.9 & 28.7 \\
%          & Ours & 7.2 & 32.5 \\ \midrule
%         Phi4-mini & Vanilla & 5.3 & 24.4 \\
%          & Ours & 12.1 & 31.3 \\ \midrule
%         Llama3-8B & Vanilla & 1.7 & 23.4 \\
%          & Ours & 5.3 & 25.8 \\ \bottomrule
% \end{tabular}
% }
% \end{wraptable} 

 
% \begin{table}[h!]
% \begin{table}[t!]
%   \caption{\method outperforms the vanilla MAS using NL and Autoform using fixed formats prompted by the LLM.} 
%     \label{tab:main_results}
%     \centering
%     % \begin{tabular}{c{0.2cm} c{0.2cm} c{0.2cm} c{0.2cm} c{0.2cm} c{0.2cm} c{0.2cm} c{0.2cm} c{0.2cm}}
%     % \resizebox{\columnwidth}{!}{
%     \resizebox{0.5\textwidth}{!}{
%     \begin{tabular}{c c c c c c c}
%       % \textbf{Base LLM} & \textbf{Method} & \textbf{GSM+} & \textbf{WikiHop} & \textbf{HotpotQA} & \textbf{NarrativeQA} & \textbf{Average} \\
%       \textbf{LLM} & \textbf{Method} & \textbf{GSM} & \textbf{WQA} & \textbf{HQA} & \textbf{NQA} & \textbf{Avg} \\
%         \hline \hline
%         % GPT-4o & Vanilla MAS & 86.5 & 10.2 & 60.7 & 47.1 & 51.125 \\
%         % & Autoform & 87.5 & 11.5 & 62.5 & 48.6 & 52.525 \\
%         % & \method (Ours) & 89.5 & 13.9 & 64.5 & 51.1 & 54.75 \\
%         % \hline
%         % o3 & Vanilla MAS & 93.5 & 15.5 & 65.2 & 53.6 & 56.95 \\
%         % & Autoform & 93.9 & 16.3 & 66.1 & 55.8 & 58.02 \\
%         % & \method (Ours) & 94.7 & 17.2 & 67.3 & 56.7 & 58.98 \\
%         % \hline
%         % Phi4-mini & Vanilla MAS & 49.2& 5.5& 29.6& 22.9& 26.8 \\
%         % & Autoform & 50.7& 6.2& 31.4& 23.8& 28.02 \\
%         % & \method (Ours) & 53.5& 8.3& 34.2& 25.9& 30.48 \\
%         % \hline   
%         GPT-4o & Vanilla & 86.5 & 19.5 & 65.1 & 51.3 & 55.6 \\
%         & Autoform & 87.5 & 18.2 & 48.5 & 35.9 & 47.5 \\
%         & \method & 89.5 & 21.4 & 69.1 & 57.2 & 59.3 \\
%         \hline
%         o3 & Vanilla & 93.5 & 20.5 & 80.2 & 70.5 & 66.2 \\
%         & Autoform & 92.1 & 21.9 & 81.7 & 59.8 & 63.9 \\
%         & \method & 94.7 & 22.8 & 88.9 & 74.3 & 70.2 \\
%         \hline
%         Phi4- & Vanilla & 49.2 & 8.3 & 24.1 & 21.5 & 25.8 \\
%         mini& Autoform & 50.7 & 12.2 & 31.5 & 12.2 & 26.6 \\
%         & \method & 53.5 & 15.7 & 33.5 & 25.5 & 32.1 \\
%         \hline
        
%         % Llama-3-8B & Vanilla MAS & 93.5 & 15.5 & 65.2 & 53.6 & 56.95 \\
%         % & Autoform & 93.9 & 16.3 & 66.1 & 55.8 & 58.02 \\
%         % & \method (Ours) & 94.7 & 17.2 & 67.3 & 56.7 & 58.98 \\
%         % \hline
%         % Mistral-7B & Vanilla MAS & 93.5 & 15.5 & 65.2 & 53.6 & 56.95 \\
%         % & Autoform & 93.9 & 16.3 & 66.1 & 55.8 & 58.02 \\
%         % & \method (Ours) & 94.7 & 17.2 & 67.3 & 56.7 & 58.98 \\
%         \hline
%     \end{tabular}
%     }
% \end{table}

% \subsection{Abalations:}
% \subsubsection{Do we even need a structuriser}

\subsection{Performance with Fixed Structures (RQ2)}
In this study, we analyze if finding the optimal structure is necessary by comparing with some largely adopted fixed structures in multi-agent communication literature. 
% The following structures are used in our experiments - KQML ~\cite{labrou1999agent}
% % ~\cite{labrou1999agent,finin1995kqml,mayfield1996evaluation}, 
% % S-expressions~\cite{delatorre2025tool},
% JSON~\cite{yin2024fast}, XML~\cite{aalst2003xml}, YAML~\cite{pathway2024yaml}, Code, Tabular, Equation, List ~\cite{chen-etal-2024-beyond-natural}.  
From table \ref{tab:fixed_format_abl} we can see that among fixed methods, Pseudo-code \& List (Sequential Thinking) generally perform well due to the mathematical nature of the tasks. However, there exists no structure that performs consistently better across all datasets. This shows that the best structure is task dependent, answering RQ2. Further our method performs better than the fixed structures indicating that \method attempts to find the optimal structure for a task.
\textcolor{blue}{We also report average tokens and time per message and across datasets. We find that alternative formats (e.g. KQML) have higher efficiency than NL and \method is able to maintain high accuracy while also finding formats that have low token count and latency.}
% KQML supports knowledge sharing by defining performatives for agent negotiation, goal expression, and flexible run-time coordination in distributed, heterogeneous multi-agent environments~\cite{labrou1999agent,finin1995kqml,mayfield1996evaluation}. Recently introduced S-expressions enable concise, compositional exchange of symbolic concepts, allowing agents to transmit structured data and logic with minimal syntactic overhead~\cite{delatorre2025tool}. JSON is widely adopted for efficient, schema-compliant data interchange between systems, minimizing message size and maximizing compatibility~\cite{yin2024fast,shorten2024structured}.XML excels at conveying complex, hierarchical information, supporting robust validation and interoperability in enterprise and workflow-heavy contexts~\cite{mishra2002structural,aalst2003xml}. 
% YAML improves human-editability and clarity in configuration and task orchestration, facilitating readable communication patterns for agent pipeline design~\cite{pathway2024yaml}. 

% % \begin{table}[h!]
% \begin{table}[t!]
%   \caption{\method outperforms fixed representations indicating the importance of dynamic learning.} 
%     \label{tab:fixed_format_abl}
%     \centering
%     % \begin{tabular}{c{0.2cm} c{0.2cm} c{0.2cm} c{0.2cm} c{0.2cm} c{0.2cm} c{0.2cm} c{0.2cm} c{0.2cm}}
%     % \resizebox{\columnwidth}{!}{
%     \resizebox{0.4\textwidth}{!}{
%     \begin{tabular}{c c c c c c c}
%       % \textbf{Structure} & \textbf{GSM+} & \textbf{WikiHop} & \textbf{HotpotQA} & \textbf{NarrativeQA} & \textbf{Average} \\
%       \textbf{Structure} & \textbf{GSM} & \textbf{WQA} & \textbf{HQA} & \textbf{NQA} & \textbf{Avg} \\
%         \hline \hline
%         % NL & 86.5 & 10.2 & 60.7 & 47.1 & 51.125 \\
%         % \method & 89.5 & 13.9 & 64.5 & 51.1 & 54.75 \\
%         % \hline
%         % KQML & 87.5 & 11.5 & 62.5 & 48.6 & 52.525 \\
%         % JSON & 93.5 & 15.5 & 65.2 & 53.6 & 56.95 \\
%         % YAML & 93.5 & 15.5 & 65.2 & 53.6 & 56.95 \\
%         % XML & 93.5 & 15.5 & 65.2 & 53.6 & 56.95 \\
%         % Code & 93.5 & 15.5 & 65.2 & 53.6 & 56.95 \\
%         % Tabular & 93.5 & 15.5 & 65.2 & 53.6 & 56.95 \\
%         % Exp & 93.5 & 15.5 & 65.2 & 53.6 & 56.95 \\
%         % List & 93.5 & 15.5 & 65.2 & 53.6 & 56.95 \\
%         % \hline
%         NL & 86.5 & 19.5 & 65.1 & 51.3 & 55.6 \\
%         Ours & \textbf{89.5} & \textbf{21.4} & \textbf{69.1} & \textbf{57.2} & \textbf{59.3} \\
%         \hline
%         JSON & 71.5 & 19.5 & 66.9 & 52.9 & 52.7 \\
%         XML & 89.1 & 17.1 & 68.1 & 52.2 & 56.6 \\
%         KQML & \textbf{89.5} & 14.6 & 63.9 & \underline{57.1} & 56.3 \\
%         YML & 87.6 & 18.3 & 66.8 & 50.0 & 55.7 \\
%         Code & \underline{89.2} & \underline{21.1} & 68.0 & 54.0 & 58.1 \\
%         Equation & 87.9 & 18.3 & \underline{69.0} & 47.4 & 55.7 \\
%         Table  & 87.0 & 14.6 & 62.2 & 53.1 & 54.2 \\
%         List & 88.1 & 20.8 & 67.9 & 56.3 & \underline{58.3} \\
%         \hline
        
%         \hline
%     \end{tabular}
%     }
% \end{table}

% \begin{table}[h!]
\begin{table}[t!]
  \caption{\method outperforms fixed representations indicating the importance of dynamic learning.} 
    \label{tab:fixed_format_abl}
    \centering
    % \begin{tabular}{c{0.2cm} c{0.2cm} c{0.2cm} c{0.2cm} c{0.2cm} c{0.2cm} c{0.2cm} c{0.2cm} c{0.2cm}}
    % \resizebox{\columnwidth}{!}{
    \resizebox{0.5\textwidth}{!}{
    \begin{tabular}{c c c c c c | c c}
      % \textbf{Structure} & \textbf{GSM+} & \textbf{WikiHop} & \textbf{HotpotQA} & \textbf{NarrativeQA} & \textbf{Average} \\
      \textbf{Structure} & \textbf{GSM} & \textbf{WQA} & \textbf{HQA} & \textbf{NQA} & \textbf{Avg} & \textbf{\textcolor{blue}{Avg Tokens}} & \textbf{\textcolor{blue}{Avg Time (s)}} \\
        \hline \hline
        % NL & 86.5 & 10.2 & 60.7 & 47.1 & 51.125 \\
        % \method & 89.5 & 13.9 & 64.5 & 51.1 & 54.75 \\
        % \hline
        % KQML & 87.5 & 11.5 & 62.5 & 48.6 & 52.525 \\
        % JSON & 93.5 & 15.5 & 65.2 & 53.6 & 56.95 \\
        % YAML & 93.5 & 15.5 & 65.2 & 53.6 & 56.95 \\
        % XML & 93.5 & 15.5 & 65.2 & 53.6 & 56.95 \\
        % Code & 93.5 & 15.5 & 65.2 & 53.6 & 56.95 \\
        % Tabular & 93.5 & 15.5 & 65.2 & 53.6 & 56.95 \\
        % Exp & 93.5 & 15.5 & 65.2 & 53.6 & 56.95 \\
        % List & 93.5 & 15.5 & 65.2 & 53.6 & 56.95 \\
        % \hline
        NL & 86.5 & 19.5 & 65.1 & 51.3 & 55.6 & 87.8 & 1.58 \\
        Ours & \textbf{89.5} & \textbf{21.4} & \textbf{69.1} & \textbf{57.2} & \textbf{59.3} & 79.8 & 1.23 \\
        \hline
        JSON & 71.5 & 19.5 & 66.9 & 52.9 & 52.7 & 101 & 1.5 \\
        XML & 89.1 & 17.1 & 68.1 & 52.2 & 56.6 & 99.8 & 1.48 \\
        KQML & \textbf{89.5} & 14.6 & 63.9 & \underline{57.1} & 56.3 & 70.5 & 1.28 \\
        YML & 87.6 & 18.3 & 66.8 & 50.0 & 55.7 & 77.3 & 1.20 \\
        Code & \underline{89.2} & \underline{21.1} & 68.0 & 54.0 & 58.1 & 201.7 & 2.41 \\
        Equation & 87.9 & 18.3 & \underline{69.0} & 47.4 & 55.7 & 88.8 & 1.58 \\
        Table  & 87.0 & 14.6 & 62.2 & 53.1 & 54.2 & 117.3 & 1.54 \\
        List & 88.1 & 20.8 & 67.9 & 56.3 & \underline{58.3} & 106.1 & 1.34 \\
        \hline
        
        \hline
    \end{tabular}
    }
\end{table}

\subsection{Token Efficiency (RQ3)}
% In this study we observe how efficient our method is compared to the vanilla multi-agent system without optimizing the message structures. This would help inform practitioners of the tradeoffs in tokens/cost with enhanced performance. 
We observe from table \ref{appendix_tab:efficiency_study} (\textcolor{blue}{Appendix}) that from all the datasets experimented, the maximum increase in token usage is not more than 20\%. In some cases \method is more efficient compared to the vanilla multi-agent system. This is, since certain formal structures require lesser tokens compared to natural language. Thus \method enhances performance while keeping the cost bounded.
% , answering RQ3.

% \begin{table}[h!]
%   \caption{Token efficiency of \method.} 
%     \label{tab:efficiency_study}
%     \centering
%     % \begin{tabular}{c{0.2cm} c{0.2cm} c{0.2cm} c{0.2cm} c{0.2cm} c{0.2cm} c{0.2cm} c{0.2cm} c{0.2cm}}
%     % \resizebox{\columnwidth}{!}{
%     \resizebox{0.45\textwidth}{!}{
%     \begin{tabular}{c c c c c}
%       \textbf{Dataset} & \textbf{\#Tokens (Vanilla)} & \textbf{\#Tokens (Ours)} & \textbf{$\Delta$ Tokens (\%)} \\
%         \hline \hline
%         GSM+ & 2667 & 2434 & -8.7  \\
%         WikiHop & 1023 & 968 & -5.4  \\
%         HotPotQA & 671 & 545 & -18.8  \\
%         NarrativeQA & 925 & 1104 & +19.3  \\
%         \hline
%     \end{tabular}
%     }
% \end{table}

\subsection{Ablation Study (RQ4)}
We study the impact of fixed set of structures/tasks (w/o dynamic) and dynamic but not learning the policy from the data (w/o learning). 
% We further study how the performance evolves with converging Q values as the learning progresses.
From figure \ref{fig:ablation} we find that the results drop without the policy learning ($\sim 42 \%$) or dynamically evolving modules. 
% On the WikiHop dataset we find the drop is drastic ($\sim 42 \%$) without the policy learning module whereas in the HotpotQA dataset we find that dynamically learning structures is more important to the task. Thus the importance of the modules are task/dataset dependent but nevertheless both the modules are important in \method to achieve optimal performance. 
From the Q-plot we can see that the Q-values/accuracy converge with time. Further, the average slope ($\Delta Q/\text{step}$) in the expansion phase is $3.08\times10^{-2}$ vs $1.5\times10^{-3}$ in the convergence phase, indicating convergence once the EMDP is frozen (supporting prop. \ref{prop_emdp_convergence}).
\textcolor{blue}{Further studies are provided in the Appendix (\ref{ablation_appendix}, \ref{ablation}).}
% In this section we perform an ablation study to analyze the importance of each component of \method. Specifically we study the impact of having a fixed set of structures/tasks and keeping the structures/tasks dynamic but not learning the policy from the environment/data. We further study how the performance evolves with converging Q values as the learning progresses.
% From figure \ref{fig:ablation} we find that the results drop without the policy learning or dynamically evolving modules. On the WikiHop dataset we find the drop is drastic ($\sim 42 \%$) without the policy learning module whereas in the HotpotQA dataset we find that dynamically learning structures is more important to the task. Thus the importance of the modules are task/dataset dependent but nevertheless both the modules are important in \method to achieve optimal performance. 
% From the Q-plot we can see that the Q-values/accuracy converge with time.
% From the Q-plot we can see that the Q values and accuracy converges as the learning progresses.

% \begin{figure}[h]
%     \centering
%      \begin{subfigure}[b]{0.22\textwidth}
%      \centering
% % \includegraphics[width=\linewidth]{figures/Cloudservices_monitors.jpg}
% % \includegraphics[width=\linewidth]{AAAI/figures/sac_ablation_old2.png}
% \includegraphics[width=\linewidth]{AAAI/figures/sac_ablation.png}
%         \caption{}
%            \label{subfig:ablation_study}
%     \end{subfigure}%
%     \begin{subfigure}[b]{0.27\textwidth}
%          \centering
% \includegraphics[width=\linewidth]{AAAI/figures/hotpotqa_qvalues.png}
%     \caption{}
%            \label{subfig:qvalues}
%     \end{subfigure}
% \caption{a) Dynamic representation evolving and policy learning module are important for performance b) The Q-values/accuracy converge as learning progresses.}
% \label{fig:ablation}
% \end{figure}
\begin{figure}[h]
    \centering
     \begin{subfigure}[b]{0.42\textwidth}
     \centering
% \includegraphics[width=\linewidth]{figures/Cloudservices_monitors.jpg}
% \includegraphics[width=\linewidth]{AAAI/figures/sac_ablation_old2.png}
\includegraphics[width=\linewidth]{AAAI/figures/sac_ablation.png}
        \caption{}
           \label{subfig:ablation_study}
    \end{subfigure}%
    
    \begin{subfigure}[b]{0.42\textwidth}
         \centering
\includegraphics[width=\linewidth]{AAAI/figures/hotpotqa_qvalues.png}
    \caption{}
           \label{subfig:qvalues}
    \end{subfigure}
\caption{a) Dynamic representation evolving and policy learning module are important for performance b) The Q-values/accuracy converge as learning progresses.}
\label{fig:ablation}
\end{figure}

\subsection{Structure Distribution (RQ5)}
% In this section we study the distribution of the tasks and structures uncovered by \method during the expansion phase of learning. 
In figure \ref{fig:structure_distribution} we see the structure distribution where we display 10(/64) high level structures including seed and discovered structures. 
% In total the method discovered 64 structures and we club the fine grained structures into these high level types. 
In addition to the seeded structures of JSON, XML, YAML, List, Code etc., our method is able to discover some new structures ($\sim 50\%$) within the categories of Trees, Graphs, Math etc. 
% The newly discovered structures account for more than 50\% of the structures. 
% We also analyze the structure distribution per task. The major tasks in the concerned dataset (HotPotQA) were mainly identified as ``Comparison'' (mathematical comparison between entities) and ``Bridge'' (multi hop reasoning between connected entities) with easy (E), medium (M) and hard categories (H). We see different tasks have different distributions of the optimal structures. 
From the task wise distribution of representations, we see the hard category of the comparison task benefits from ``Math'' representation due to the nature of the task, whereas the bridge task benefits from Graph and Tree structures as expected. Thus \method is able to select the pertinent structure for the task. 
% In this section we study the distribution of the tasks and structures uncovered by \method during the expansion phase of learning. In figure \ref{fig:structure_distribution} we see the structure distribution where we display 10 high level structures including seed and discovered structures. In total the method discovered 64 structures and we club the fine grained structures into these high level types. In addition to the seeded structures of JSON, XML, YAML, List, Code etc., our method is able to discover some new structures within the categories of Trees, Graphs, Math etc. The newly discovered structures account for more than 50\% of the structures. We also analyze the structure distribution per task. The major tasks in the concerned dataset (HotPotQA) were mainly identified as ``Comparison'' (mathematical comparison between entities) and ``Bridge'' (multi hop reasoning between connected entities) with easy (E), medium (M) and hard categories (H). We see different tasks have different distributions of the optimal structures. The hard category of the comparison task benefits from ``Math'' representation due to the nature of the task, whereas the bridge task benefits from Graph and Tree structures as expected. Thus \method is able to select the pertinent structure for the task. 
% For brevity, we add further details in the \textcolor{blue}{Appendix \ref{apndx_format_evolution}}.
\textcolor{blue}{For brevity, we show evolution of formats in the \textcolor{blue}{Appendix \ref{apndx_format_evolution}}.}
% Due to space constraints, we add further experiment results \& details in the appendix.
% \footnote{Appendix is submitted as a supplementary to the main paper.}.

\begin{figure}[h]
    \centering
     \begin{subfigure}[b]{0.24\textwidth}
     \centering
% \includegraphics[width=\linewidth]{figures/Cloudservices_monitors.jpg}
\includegraphics[width=\linewidth]{AAAI/figures/sac_structure.png}
        \caption{}
           \label{subfig:structure_dist}
    \end{subfigure}%
    \begin{subfigure}[b]{0.25\textwidth}
         \centering
\includegraphics[width=\linewidth]{AAAI/figures/sac_structure_for_task.png}
    \caption{}
           \label{subfig:struct_for_task}
    \end{subfigure}
% \caption{Fig. a) shows the overall structure distribution found by \method for HotPotQA dataset. Fig. b) shows the structure distribution by task type.}
% \caption{Fig. a) shows the overall structure distribution and b)is the task-wise distribution found by \method for HotPotQA dataset.}
\caption{Fig. a) shows the overall structure distribution and b) is the task-wise distribution found for HotPotQA.}
\label{fig:structure_distribution}
\end{figure}

% \begin{figure}[h]
%     \centering
%      \begin{subfigure}[b]{0.3\textwidth}
%      \centering
% % \includegraphics[width=\linewidth]{figures/Cloudservices_monitors.jpg}
% \includegraphics[width=\linewidth]{AAAI/figures/sac_structure.png}
%         \caption{}
%            \label{subfig:structure_dist}
%     \end{subfigure}%
    
%     \begin{subfigure}[b]{0.45\textwidth}
%          \centering
% \includegraphics[width=\linewidth]{AAAI/figures/sac_structure_for_task.png}
%     \caption{}
%            \label{subfig:struct_for_task}
%     \end{subfigure}
% % \caption{Fig. a) shows the overall structure distribution found by \method for HotPotQA dataset. Fig. b) shows the structure distribution by task type.}
% \caption{Fig. a) shows the overall structure distribution and b)is the task-wise distribution found for HotPotQA dataset.}
% \label{fig:structure_distribution}
% \end{figure}

% \subsubsection{Implicit vs Explicit Struturisation} In our experiments, we first generate the natural language response and then convert it to a suitable structure based on the suggested format. This was done so that we could compare various metrics [ADD METRICS] that quantify structures in comparison with Natural-Language. We run these abalsations demonstrating similar performance when we implicitly embed the suggested structure in the prompt and directly generate the structured response like Autoform, with the difference being that. The reason why both are close is that LLMs use same inductive biases about the recipe to convert to suggested structure, same system-prompt, and similar message histories. 

% [ADD RESULTS HERE]

% \subsubsection{Use of Different Models}

% In this ablation, we explore the role of different model scales—Small Language Models (SLMs) and Large Language Models (LLMs)—across three subtasks: (i) task prediction (states), (ii) exploration of new formats, and (iii) applying selected formats (actions). Empirical observations and previous work motivate three key hypotheses.

% \begin{enumerate}
% \item \textbf{SLMs are effective for domain adaptation and offer broader coverage of search space during exploration}, due to their higher variance and local adaptability~\cite{marquez2025fine,mwase2023uncovering,mojarradi2024improving}. However, they tend to struggle with out-of-distribution (OOD) generalization and exhibit inconsistencies, particularly in knowledge-intensive settings~\cite{kang2023knowledge,li2024simple}.

% text
% \item \textbf{LLMs are more robust in generating consistent actions and performing under distribution shift}, benefiting from instruction tuning and larger pretraining corpora~\cite{sun2023evaluating,agrawal2025enhancing,howe2024effects}. As a result, they demonstrate less sensitivity to prompt-level structurization, limiting its marginal utility~\cite{wang2025slot,geng2023grammar}.

% \item \textbf{The effectiveness of optimal structured prompts depends highly on the training mode and architectural differences between model families}. Bidirectional encoder models like BERT and its derivatives, trained on masked language modeling objectives, require fundamentally different prompting strategies compared to autoregressive models like GPT~\cite{patel2022bidirectional,webson2022prompt}. BERT-based models excel with cloze-style prompts and mask-based templates that leverage their bidirectional context understanding~\cite{jiang2022promptbert,qi2023multi}, while GPT models perform optimally with sequential, instruction-following prompts that align with their left-to-right generation paradigm~\cite{brown2020language,wei2022chain}. This architectural sensitivity extends to prompt structure optimization, where encoder-only models benefit from structured template designs~\cite{liu2023structured,sun2022how}, whereas decoder-only models show greater flexibility with natural language instructions but require careful prompt engineering to avoid hallucinations~\cite{schulhoff2024prompt,li2023exploring}.
% \end{enumerate}

% Suggested Models for Evaluation:

% For comprehensive evaluation across different architectural paradigms, we recommend testing the following models:

% \textbf{Autoregressive (Decoder-only) Models:}
% \begin{itemize}
% \item \textbf{GPT-4o}: Advanced multimodal capabilities with enhanced instruction following~\cite{openai2024gpt4o}
% \item \textbf{GPT-3.5-turbo}: Cost-effective baseline for autoregressive prompting~\cite{ouyang2022training}
% \item \textbf{Llama-3 (8B/70B)}: Open-source alternative with strong reasoning capabilities~\cite{touvron2023llama}
% \item \textbf{Phi-4}: Microsoft's efficient reasoning model with structured thinking capabilities~\cite{microsoft2024phi4}
% \end{itemize}

% \textbf{Encoder-only Models:}
% \begin{itemize}
% \item \textbf{BERT-large}: Bidirectional baseline for masked language modeling approaches~\cite{devlin2019bert}
% \item \textbf{RoBERTa-large}: Optimized BERT variant with improved pretraining~\cite{liu2019roberta}
% \end{itemize}

% \textbf{Encoder-Decoder Models:}
% \begin{itemize}
% \item \textbf{Gemini-Pro}: Google's multimodal model with bidirectional reasoning~\cite{team2023gemini}
% \item \textbf{T5-large}: Text-to-text unified framework for structured tasks~\cite{raffel2020exploring}
% \end{itemize}

% [ADD RESULTS HERE]
\section{Conclusion}
In this paper, we present a method to dynamically learn the optimal structure for effective communication in a multi-agent system. We pose the problem of finding the optimal message structure as an Expanding Markov Decision Process (EMDP). This allows us to dynamically evolve to new representations as the agents interact with the environment. We further learn the optimal policy to decide the structure for the task at hand, rather than relying on fixed structures or the inductive biases from the LLM. Ours is the first work to provide a learnable framework for optimizing message representations in multi-agent (LLM) communication. 
We plan to extend this research by employing composite optimal representations and latent search spaces to expand beyond finite tabular methods to handle infinite tasks and representations.
% We plan to extend this research to employ composite optimal structures. Future research could also explore the latent search space for expanding to infinite tasks and structures beyond the finite set using tabular methods.

% In this paper, we present a method to dynamically learn the optimal structure for effective communication in a multi-agent system. The formulation of the problem of finding the optimal message structure as an Expanding Markov Decision Process (EMDP) allows us to dynamically evolve to new structures as the agents interact with the environment. We further learn the optimal policy to decide the structure for the task at hand, rather than relying on fixed structures or the inductive biases from the LLM. Ours is the first work to provide a learnable framework for optimizing message representations in multi-agent (LLM) communication. 
% We plan to extend this research by employing composite optimal structures and latent search spaces to expand beyond finite tabular methods to handle infinite tasks and structures.
% % We plan to extend this research to employ composite optimal structures. Future research could also explore the latent search space for expanding to infinite tasks and structures beyond the finite set using tabular methods.


\bibliography{references_arr}

\appendix


\end{document}
